\documentclass[conference,compsoc]{IEEEtran}
%DIF LATEXDIFF DIFFERENCE FILE
%DIF DEL old\Paper\main.tex   Fri Aug 21 11:47:08 2026
%DIF ADD new\Paper\main.tex   Fri Aug 21 11:27:28 2026

\usepackage[T1]{fontenc}
\usepackage{graphicx}
\usepackage{cite}
\usepackage[n,advantage,operators,sets,
  adversary,landau,probability,notions,
  logic,ff,mm,primitives,events,complexity,
  oracles,asymptotics,keys]{cryptocode}
\usepackage{amsfonts}   
\usepackage{mathtools} 
\usepackage{amsthm}
\usepackage{booktabs}
\usepackage{algorithm}
\usepackage{algpseudocode}
\usepackage{placeins}
\usepackage{tcolorbox}
\usepackage{tabularx}
\usepackage{adjustbox}
\usepackage[inline]{enumitem}
\usepackage{multirow}
\usepackage{pgfplots}
\pgfplotsset{compat=1.18}
\usepgfplotslibrary{groupplots}
\usepackage[hidelinks]{hyperref}
\usepackage{float}
\usepackage{amsmath}

% 공통으로 사용할 예쁜 블록 스타일 정의
\newtcbox{\myblock}[1][]{
  colback=blue!5!white,    % 배경색 (연한 파란색)
  colframe=blue!60!black,  % 테두리색
  fontupper=\sffamily,     % 깔끔한 고딕(Sans-serif) 폰트
  nobeforeafter, 
  center, 
  tcbox raise base,
  arc=4pt,                 % 둥근 모서리 반지름
  boxrule=1pt,             % 테두리 두께
  top=4pt, bottom=4pt, left=8pt, right=8pt % 내부 여백
}
%% ── Commitment & Merkle ──────────────────────────────────────────────────────
\newcommand{\para}{\textsf{pp}}
\newcommand{\cm}{\textsf{cm}}
\newcommand{\rt}{\textsf{rt}}
\newcommand{\cmmt}{\textsf{cm}_{\textsf{mt}}}
\newcommand{\merkleC}{\textsf{MT.Com}}
\newcommand{\merkleV}{\textsf{MT.Verify}}
\newcommand{\merkleO}{\textsf{MT.Open}}
\newcommand{\pedersen}{\textsf{P.CM}}
\newcommand{\ck}{\textsf{ck}}
\newcommand{\pedC}{\textsf{Ped.Com}}
\newcommand{\pedO}{\textsf{Ped.Open}}
\newcommand{\pedV}{\textsf{Ped.Verify}}

%% ── Algebra ──────────────────────────────────────────────────────────────────
\newcommand{\mat}{\mathbf}          % matrix: \mat{A}
\newcommand{\vect}{\mathbf}         % vector: \vect{x}
%% \enc is already defined by cryptocode's 'primitives' module (as encryption).
%% We redefine it here to mean the linear-code encoding operator instead.
\renewcommand{\enc}{\mathsf{Enc}}   % encoding operator (overrides cryptocode)

%% ── SNARK ────────────────────────────────────────────────────────────────────
\newcommand{\snarkP}{\textsf{SNARK.P}}
\newcommand{\snarkV}{\textsf{SNARK.V}}
\newcommand{\bPi}{{\bf \Pi}}
\newcommand{\setup}{{\sf Setup}}
\newcommand{\com}{{\sf Com}}
\newcommand{\prove}{{\sf Prove}}
\renewcommand{\verify}{{\sf Verify}}
\newcommand{\relation}{{\mathcal{R}}}
\newcommand{\cp}{{\mathsf{cp}}}
\renewcommand{\pk}{{\textsf{pk}}}
\renewcommand{\vk}{{\textsf{vk}}}
\renewcommand{\ck}{{\textsf{ck}}}



%% ── SNARK statement / witness (replaces \mathbb{x}, \mathbb{w}) ─────────────
%% \mathbb works only for uppercase.
%% We use \mathsf to distinguish SNARK instances from vector \vect{x} (\mathbf{x}).
\newcommand{\stmt}{\mathsf{x}}      % SNARK public statement
\newcommand{\wit}{\mathsf{w}}       % SNARK witness
\newcommand{\comwit}{\mathsf{u}}

%% ── Misc ─────────────────────────────────────────────────────────────────────
\newcommand{\fold}{\textsf{Fold}}
\newcommand{\Hash}{\textsf{Hash}}
\newcommand{\vecBindV}{\textsf{Vec.Bind}}
\newcommand{\cpLinkP}{\textsf{CPLink.Prove}}
\newcommand{\cpLinkV}{\textsf{CPLink.Verify}}
\newcommand{\protocol}{\textsf{LAMP}}   % protocol name
\newcommand{\fmMerkle}{\protocol_{\mathsf{Merkle}}}
\newcommand{\fmKZG}{\protocol_{\mathsf{KZG}}}
\newcommand{\fmCPLink}{\protocol_{\mathsf{CPLink}}}
\newcommand{\fmQA}{\protocol_{\mathsf{QA}}}

%% ── Sets of objects (avoid \mathcal{V} collision with verifier) ──────────────
%% Vector set is written explicitly in text to avoid notation collision.

%% ── Theorem environments ─────────────────────────────────────────────────────
\theoremstyle{plain}
\newtheorem{theorem}{Theorem}[section]
\newtheorem{lemma}[theorem]{Lemma}
\newtheorem{claim}[theorem]{Claim}
\newtheorem{corollary}[theorem]{Corollary}

\theoremstyle{definition}
\newtheorem{definition}[theorem]{Definition}

\theoremstyle{remark}
\newtheorem{remark}{Remark}

%% Some conference templates provide \Description for accessibility metadata.
\providecommand{\Description}[1]{}
%DIF PREAMBLE EXTENSION ADDED BY LATEXDIFF
%DIF CFONT PREAMBLE %DIF PREAMBLE
\RequirePackage{color}\definecolor{RED}{rgb}{1,0,0}\definecolor{DIFADD}{RGB}{17,34,51} %DIF PREAMBLE
\DeclareOldFontCommand{\sf}{\normalfont\sffamily}{\mathsf} %DIF PREAMBLE
\providecommand{\DIFaddtex}[1]{{\protect\color{DIFADD} \sf #1}} %DIF PREAMBLE
\providecommand{\DIFdeltex}[1]{{\protect\color{red} \scriptsize #1}} %DIF PREAMBLE
%DIF SAFE PREAMBLE %DIF PREAMBLE
\providecommand{\DIFaddbegin}{} %DIF PREAMBLE
\providecommand{\DIFaddend}{} %DIF PREAMBLE
\providecommand{\DIFdelbegin}{} %DIF PREAMBLE
\providecommand{\DIFdelend}{} %DIF PREAMBLE
\providecommand{\DIFmodbegin}{} %DIF PREAMBLE
\providecommand{\DIFmodend}{} %DIF PREAMBLE
%DIF FLOATSAFE PREAMBLE %DIF PREAMBLE
\providecommand{\DIFaddFL}[1]{\DIFadd{#1}} %DIF PREAMBLE
\providecommand{\DIFdelFL}[1]{\DIFdel{#1}} %DIF PREAMBLE
\providecommand{\DIFaddbeginFL}{} %DIF PREAMBLE
\providecommand{\DIFaddendFL}{} %DIF PREAMBLE
\providecommand{\DIFdelbeginFL}{} %DIF PREAMBLE
\providecommand{\DIFdelendFL}{} %DIF PREAMBLE
%DIF HYPERREF PREAMBLE %DIF PREAMBLE
\providecommand{\DIFadd}[1]{\texorpdfstring{\DIFaddtex{#1}}{#1}} %DIF PREAMBLE
\providecommand{\DIFdel}[1]{\texorpdfstring{\DIFdeltex{#1}}{}} %DIF PREAMBLE
\providecommand{\DIFscaledelfig}{0.5}
%DIF HIGHLIGHTGRAPHICS PREAMBLE %DIF PREAMBLE
\RequirePackage{settobox} %DIF PREAMBLE
\RequirePackage{letltxmacro} %DIF PREAMBLE
\newsavebox{\DIFdelgraphicsbox} %DIF PREAMBLE
\newlength{\DIFdelgraphicswidth} %DIF PREAMBLE
\newlength{\DIFdelgraphicsheight} %DIF PREAMBLE
% store original definition of \includegraphics %DIF PREAMBLE
\LetLtxMacro{\DIFOincludegraphics}{\includegraphics} %DIF PREAMBLE
\providecommand{\DIFaddincludegraphics}[2][]{{\color{blue}\fbox{\DIFOincludegraphics[#1]{#2}}}} %DIF PREAMBLE
\providecommand{\DIFdelincludegraphics}[2][]{% %DIF PREAMBLE
\sbox{\DIFdelgraphicsbox}{\DIFOincludegraphics[#1]{#2}}% %DIF PREAMBLE
\settoboxwidth{\DIFdelgraphicswidth}{\DIFdelgraphicsbox} %DIF PREAMBLE
\settoboxtotalheight{\DIFdelgraphicsheight}{\DIFdelgraphicsbox} %DIF PREAMBLE
\scalebox{\DIFscaledelfig}{% %DIF PREAMBLE
\parbox[b]{\DIFdelgraphicswidth}{\usebox{\DIFdelgraphicsbox}\\[-\baselineskip] \rule{\DIFdelgraphicswidth}{0em}}\llap{\resizebox{\DIFdelgraphicswidth}{\DIFdelgraphicsheight}{% %DIF PREAMBLE
\setlength{\unitlength}{\DIFdelgraphicswidth}% %DIF PREAMBLE
\begin{picture}(1,1)% %DIF PREAMBLE
\thicklines\linethickness{2pt} %DIF PREAMBLE
{\color[rgb]{1,0,0}\put(0,0){\framebox(1,1){}}}% %DIF PREAMBLE
{\color[rgb]{1,0,0}\put(0,0){\line( 1,1){1}}}% %DIF PREAMBLE
{\color[rgb]{1,0,0}\put(0,1){\line(1,-1){1}}}% %DIF PREAMBLE
\end{picture}% %DIF PREAMBLE
}\hspace*{3pt}}} %DIF PREAMBLE
} %DIF PREAMBLE
\LetLtxMacro{\DIFOaddbegin}{\DIFaddbegin} %DIF PREAMBLE
\LetLtxMacro{\DIFOaddend}{\DIFaddend} %DIF PREAMBLE
\LetLtxMacro{\DIFOdelbegin}{\DIFdelbegin} %DIF PREAMBLE
\LetLtxMacro{\DIFOdelend}{\DIFdelend} %DIF PREAMBLE
\DeclareRobustCommand{\DIFaddbegin}{\DIFOaddbegin \let\includegraphics\DIFaddincludegraphics} %DIF PREAMBLE
\DeclareRobustCommand{\DIFaddend}{\DIFOaddend \let\includegraphics\DIFOincludegraphics} %DIF PREAMBLE
\DeclareRobustCommand{\DIFdelbegin}{\DIFOdelbegin \let\includegraphics\DIFdelincludegraphics} %DIF PREAMBLE
\DeclareRobustCommand{\DIFdelend}{\DIFOaddend \let\includegraphics\DIFOincludegraphics} %DIF PREAMBLE
\LetLtxMacro{\DIFOaddbeginFL}{\DIFaddbeginFL} %DIF PREAMBLE
\LetLtxMacro{\DIFOaddendFL}{\DIFaddendFL} %DIF PREAMBLE
\LetLtxMacro{\DIFOdelbeginFL}{\DIFdelbeginFL} %DIF PREAMBLE
\LetLtxMacro{\DIFOdelendFL}{\DIFdelendFL} %DIF PREAMBLE
\DeclareRobustCommand{\DIFaddbeginFL}{\DIFOaddbeginFL \let\includegraphics\DIFaddincludegraphics} %DIF PREAMBLE
\DeclareRobustCommand{\DIFaddendFL}{\DIFOaddendFL \let\includegraphics\DIFOincludegraphics} %DIF PREAMBLE
\DeclareRobustCommand{\DIFdelbeginFL}{\DIFOdelbeginFL \let\includegraphics\DIFdelincludegraphics} %DIF PREAMBLE
\DeclareRobustCommand{\DIFdelendFL}{\DIFOaddendFL \let\includegraphics\DIFOincludegraphics} %DIF PREAMBLE
%DIF COLORLISTINGS PREAMBLE %DIF PREAMBLE
\RequirePackage{listings} %DIF PREAMBLE
\RequirePackage{color} %DIF PREAMBLE
\lstdefinelanguage{DIFcode}{ %DIF PREAMBLE
%DIF DIFCODE_CFONT %DIF PREAMBLE
  moredelim=[il][\color{red}\scriptsize]{\%DIF\ <\ }, %DIF PREAMBLE
  moredelim=[il][\color{blue}\sffamily]{\%DIF\ >\ } %DIF PREAMBLE
} %DIF PREAMBLE
\lstdefinestyle{DIFverbatimstyle}{ %DIF PREAMBLE
	language=DIFcode, %DIF PREAMBLE
	basicstyle=\ttfamily, %DIF PREAMBLE
	columns=fullflexible, %DIF PREAMBLE
	keepspaces=true %DIF PREAMBLE
} %DIF PREAMBLE
\lstnewenvironment{DIFverbatim}[1][]{\lstset{style=DIFverbatimstyle}}{} %DIF PREAMBLE
\lstnewenvironment{DIFverbatim*}[1][]{\lstset{style=DIFverbatimstyle,showspaces=true}}{} %DIF PREAMBLE
\lstset{extendedchars=\true,inputencoding=utf8}

%DIF END PREAMBLE EXTENSION ADDED BY LATEXDIFF

\begin{document}

\title{LAMP: Linear Verification of Matrix Multiplication via Proximity Testing}

\author{\IEEEauthorblockN{Anonymous Submission}}

\maketitle

\begin{abstract}
Verifiable computation systems often need to prove large matrix multiplication
statements, but a direct SNARK arithmetization of a \(k \times k\) product
requires \(\mathcal{O}(k^3)\) constraints. Freivalds' randomized check reduces
the algebraic computation to vector-matrix products, but proving those products
inside a SNARK still costs \(\mathcal{O}(k^2)\) constraints.

We present \(\protocol\), a matrix-multiplication checking protocol that combines
Freivalds' randomized check with proximity testing over linear error-correcting
codes. The prover commits to encoded matrices and intermediate vectors before
the sampled query positions are derived. The CP-SNARK circuit then checks only the sampled codeword positions and commits to the values used inside the circuit, while Merkle openings and CP-Link proofs ensure consistency between the in-circuit witnesses and the externally committed values. We prove soundness \DIFaddbegin \DIFadd{of the public-coin interactive protocol }\DIFaddend for this committed-input setting under the soundness of the SNARK backend, the binding of the commitments, the correctness of the CP-Link checks, and the distance of the code.

For a fixed number \(t\) of sampled positions, the main in-circuit SNARK
relation has  \DIFaddbegin \DIFadd{\(\mathcal{O}(tk+E_{\mathsf{code}})\) constraints, which is linear
in \(k\) for fixed \(t\) in our fixed-rate Reed--Solomon instantiation.  }\DIFaddend Merkle
openings and CP-Link  \DIFaddbegin \DIFadd{add
\(\mathcal{O}(t\log n)+E_{\mathsf{link}}(k,t)\) backend work}\DIFaddend . We implement
\(\protocol\) in Go and compare it with a Freivalds-based SNARK circuit. In the
matrix benchmark at \(k=2^{12}\), \(\protocol\) reduces the constraint count by
\(30.3\times\) and shortens proof generation time by  \DIFaddbegin \DIFadd{\(8.43\times\)}\DIFaddend ; verification
stays  \DIFaddbegin \DIFadd{around \(0.06\) }\DIFaddend seconds in the measured range.  \DIFaddbegin \DIFadd{On a GPT-2 layer,
\(\protocol\) shortens proving time by
\(1.77\times\), but verification increases from \(0.49\) s to \(6.52\) s and
proof size increases from \(196\) B to \(3.34\) MB; this workload therefore
exhibits a prover-side rather than an end-to-end improvement.
}\DIFaddend 

\end{abstract}

\begin{IEEEkeywords}
Verifiable computation, SNARKs, Proximity testing, Linear codes, Freivalds' algorithm, Merkle trees, Commitment linking
\end{IEEEkeywords}

%% Section 1: Introduction
\section{Introduction}
\label{sec:intro}

Matrix multiplication is a core computational primitive underlying many large-scale computations, including numerical linear algebra, signal processing, computer graphics, and artificial intelligence.  Consequently, many applications require verifiability for matrix products while keeping the underlying matrices private.  Zero-knowledge proofs provide a cryptographic mechanism for this goal:
a prover can convince a verifier that a statement is correct without revealing any secret information that witnesses the statement.

Applying general-purpose ZKP systems directly to large matrix multiplication, however, remains expensive.  Succinct proof systems such as ~\cite{groth2016size, plonk, maller2019sonic, chiesa2020marlin} prove relations after the computation has been represented as an arithmetic circuit or a similar algebraic constraint system.  
A standard arithmetization of a \(k\times k\) matrix product \(\mat{A}\mat{B}=\mat{C}\) computes every output
entry and therefore uses \(\mathcal{O}(k^3)\) arithmetic constraints.  This is
already impractical for moderate \(k\), and it is especially problematic for
Transformer-based neural networks~\cite{transformer, bert, GPT3}, whose projection and attention layers often
contain matrix multiplications with dimension \(k\ge 1{,}000\). \\

\noindent\textbf{The \(\mathcal{O}(k^2)\) barrier.}
Freivalds' classical randomized check~\cite{freivalds1977probabilistic} reduces
matrix-product verification from a cubic computation to a quadratic one.  Given
matrices \(\mat{A},\mat{B},\mat{C}\) and a random vector \(\vect{r}\), the
verifier checks
\[
  \vect{r}\mat{A}\mat{B}=\vect{r}\mat{C}
\]
instead of recomputing the full product.  This reduces the verification cost
from \(\mathcal{O}(k^3)\) to \(\mathcal{O}(k^2)\).  The resulting
\(\mathcal{O}(k^2)\)-size circuit is an improvement, but it is still expensive:
for example, at \(k=2^{12}\), it already requires \(50.5\) million R1CS
constraints and \(405.04\,\text{s}\) of proving time with Groth16.  Thus, a large gap remains
between the \(\mathcal{O}(k^2)\) SNARK cost and practical large-scale
matrix-multiplication verification. \\

\noindent\textbf{Our insight: from computation to checking.}
We observe that the \(\mathcal{O}(k^2)\) barrier is not inherent to the
verification problem itself.  It arises because existing SNARK-based approaches
still \emph{compute} the vector--matrix products inside the circuit.  Our key
idea is to move from ``computing inside the circuit'' to ``\emph{checking} inside the
circuit.''  The prover performs the expensive matrix arithmetic outside the
SNARK and commits to the relevant encoded data and intermediate values.  The
verifier then checks only a small number of lightweight algebraic relations
that are sufficient to detect inconsistent claims.
The technical core of this approach is proximity testing over linear
error-correcting codes.  Let \(\mathcal{C}\) be a linear code with relative
distance \(\delta\).  We encode each matrix row-wise and consider the Freivalds
intermediate vectors
\[
  \vect{x}=\vect{r}\mat{A},\qquad
  \vect{y}=\vect{x}\mat{B},\qquad
  \vect{z}=\vect{r}\mat{C}.
\]
The linearity of \(\mathcal{C}\) lets us check these vector--matrix products in
the encoded domain, column by column.  For one encoded column, a single inner
product, or \emph{fold}, checks one coordinate of the encoded product \DIFaddbegin \DIFadd{, and
random queries detect inconsistent relations according to the code distance}\DIFaddend .

The key point is that the SNARK circuit never takes the full matrices as
witnesses.  Instead, it checks fold equations on \(t\) sampled encoded columns,
where \(t\) is independent of the matrix dimension \(k\) \DIFaddbegin \DIFadd{, and checks the
encoding consistency }\DIFaddend of the intermediate vectors .
To bind the circuit witnesses to the data fixed before the query positions are known,
we use a commit-and-prove SNARK for the sampled block witnesses.  The prover
also supplies Merkle  \DIFaddbegin \DIFadd{openings and }\DIFaddend CP-Link proofs  \DIFaddbegin \DIFadd{that bind the sampled circuit
witnesses to the pre-query commitments}\DIFaddend .

% As a result, the main SNARK relation has \(\mathcal{O}(tk)\) constraints.
% For fixed \(t\), this is linear in the matrix dimension.  The sampled opening
% and linking layer adds \(\mathcal{O}(t\log n)+E_{\mathsf{link}}(k,t)\) backend
% work.  In our implementation, verification consists of \(\mathcal{O}(1)\) pairings for SNARK and CP-Link verification and \(\mathcal{O}(t\log n)\) hash operations for Merkle openings.\\

As a result, the main SNARK relation has
\(\mathcal{O}(tk+E_{\mathsf{code}})\) constraints, where
\(E_{\mathsf{code}}\) is the cost of checking that the
intermediate vectors are valid encodings under the error-correcting code. For the fixed-rate Reed--Solomon setting used in our implementation,
this remains linear in  \DIFaddbegin \DIFadd{\(k\) }\DIFaddend for fixed \(t\).   \DIFaddbegin \DIFadd{Here ``linear'' refers only to
the main in-circuit constraint count: native vector--matrix products and
encoded-matrix commitments still require \(\mathcal{O}(k^2)\) prover work.}\DIFaddend \\


\noindent\textbf{Contributions.}
We present \(\protocol\), a protocol for efficiently proving
matrix-multiplication statements.  Our contributions are as follows.
\begin{itemize}
  \item \textbf{A code-based matrix-product checking protocol.}
  \(\protocol\) combines Freivalds' randomized reduction with proximity testing
  over linear error-correcting codes.  Instead of computing vector--matrix
  products inside the SNARK circuit, the protocol checks sampled fold equations
  over encoded commitments and verifies code consistency for the intermediate
  vectors. 

  For fixed \(t\), the main circuit cost is
  \(\mathcal{O}(tk+E_{\mathsf{code}})\), which is linear in \(k\) for the
  fixed-rate Reed--Solomon setting used in our implementation.

  \item \textbf{Soundness analysis.}
  We prove soundness  \DIFaddbegin \DIFadd{of the public-coin interactive protocol for committed
  inputs, accounting for input proximity, sampled relation checks, and the
  commitment-linking backend}\DIFaddend .

  \item \textbf{Implementation and evaluation.}
  We implement \(\protocol\) in Go and compare it with a Freivalds-based SNARK
  circuit for matrix dimensions \(k=2^7\) through \(2^{13}\).  At
\(k=2^{12}\), \(\protocol\) reduces the constraint count by \(30.3\times\)
and proving time by \(8.43\times\), while verification remains independent of
the matrix dimension and stays around \(0.06\,\text{s}\).  We also evaluate
  \(\protocol\) on a GPT-2 workload with sequence length \(2^{10}\), where it
  reduces the constraint count by \(2.3\times\) and proving time by
   \DIFaddbegin \DIFadd{\(1.77\times\).  This prover-side improvement comes with a verifier-side
  trade-off: verification increases from \(0.49\) s to \(6.52\) s and proof
  size from \(196\) B to \(3.34\) MB.  Thus, the GPT-2 result demonstrates
  applicability to a }\DIFaddend matrix-heavy neural-network  \DIFaddbegin \DIFadd{workload, but not an
  end-to-end performance improvement}\DIFaddend .
\end{itemize}



 %% Section 2: Related Work
\section{Related Work}
\label{sec:related}

\noindent\textbf{Matrix multiplication.}
Several proof systems have studied how to verify matrix multiplication more
efficiently than by arithmetizing the full cubic computation. Thaler's
sum-check-based protocol~\cite{thaler2013time} shows that matrix multiplication
can be verified with quadratic prover work by reducing the claim to checks over
low-degree extensions. This avoids the naive \(\mathcal{O}(n^3)\) arithmetic
cost, but the verifier's work still grows linearly with the matrix dimension,
which is costly when the goal is succinct verification for large committed
matrices.
More recent protocols design proof systems directly for committed matrix
multiplication. zkMatrix~\cite{cong2024zkmatrix} reduces committed
 \DIFaddbegin \DIFadd{\(k\times k\) }\DIFaddend matrix multiplication to  inner-product relations and  \DIFaddbegin \DIFadd{adapts
Bulletproofs-style arguments, achieving \(\mathcal{O}(k^2)\) prover work and
\(\mathcal{O}(\log k)\) proof size and verifier work. Its follow-up,
DualMatrix~\cite{cong2026dualmatrix}, reduces the field-operation cost to
\(\mathcal{O}(K+k)\), where \(K\) is the number of nonzero matrix entries, while
using \(\mathcal{O}(k)\) group operations. For dense matrices,
\(K=\Theta(k^2)\). }\DIFaddend zkMaP~\cite{zkMaP} instead uses KZG commitments to obtain
 \DIFaddbegin \DIFadd{\(\mathcal{O}(k^2)\) prover work, a constant-size proof, and constant verifier
work without placing the matrix product in an arithmetic circuit.
}

\DIFadd{Table~\ref{tab:related-matrix-comparison} separates total prover work from
in-circuit constraints. This distinction is central to \(\protocol\): computing
the intermediate vectors and encoded-column commitments still requires
\(\mathcal{O}(k^2)\) total prover work, but the matrix-multiplication relation
inside the CP-SNARK uses only
\(\mathcal{O}(tk+E_{\mathsf{code}})\) constraints. For fixed \(t\) and
\(n=\Theta(k)\), its Merkle-dominated proof size and verifier work are
\(\mathcal{O}(\log k)\). Thus, the table does not claim an unconditional
end-to-end advantage over every prior system; it identifies the different cost
boundary improved by \(\protocol\).
}

\begin{table*}[t]
\centering
\footnotesize
\begin{adjustbox}{max width=\textwidth}
\begin{tabular}{lcccc}
\toprule
\DIFaddFL{\textbf{System} }& \DIFaddFL{\textbf{Total prover work} }& \DIFaddFL{\textbf{Main circuit constraints}
}& \DIFaddFL{\textbf{Proof size} }& \DIFaddFL{\textbf{Verifier work} }\DIFaddendFL \\
\midrule
\DIFaddFL{zkMatrix~\cite{cong2024zkmatrix}
}& \DIFaddFL{\(\mathcal{O}(k^2)\) }& \DIFaddFL{-- }& \DIFaddFL{\(\mathcal{O}(\log k)\)
}& \DIFaddFL{\(\mathcal{O}(\log k)\) }\\
\DIFaddFL{DualMatrix~\cite{cong2026dualmatrix}
}& \DIFaddFL{\(\mathcal{O}(K+k)\,\mathbb{F}+\mathcal{O}(k)\,\mathbb{G}\) }& \DIFaddFL{--
}& \DIFaddFL{\(\mathcal{O}(\log k)\) }& \DIFaddFL{\(\mathcal{O}(\log k)\) }\\
\DIFaddFL{zkMaP~\cite{zkMaP}
}& \DIFaddFL{\(\mathcal{O}(k^2)\) }& \DIFaddFL{-- }& \DIFaddFL{\(\mathcal{O}(1)\)
}& \DIFaddFL{\(\mathcal{O}(1)\) }\\
\DIFaddFL{\(\protocol\)
}& \DIFaddFL{\(\mathcal{O}(k^2)\)
}& \DIFaddFL{\(\mathcal{O}(tk+E_{\mathsf{code}})\)
}& \DIFaddFL{\(\mathcal{O}(\log k)\) for fixed \(t\)
}& \DIFaddFL{\(\mathcal{O}(\log k)\) for fixed \(t\) }\\
\bottomrule
\end{tabular}
\end{adjustbox}
\vspace{4pt}
\caption{\DIFaddFL{Asymptotic comparison with dedicated committed-matrix-multiplication
proofs. A dash means that the approach does not use a separate arithmetic
circuit for the matrix-multiplication relation. For DualMatrix,
\(K=\Theta(k^2)\) on dense matrices.}}
\label{tab:related-matrix-comparison}
\end{table*}
\DIFaddend 

\DIFaddbegin \DIFadd{We found no public zkMatrix implementation, and the zkMaP artifact link was
unavailable. We therefore use the public DualMatrix implementation as the
closest reproducible dedicated baseline in Section~\ref{sec:implementation}.
Our Freivalds-in-Groth16 baseline is not state of the art; it uses the same
Groth16 stack to isolate the effect of moving vector--matrix products out of
the circuit. }\\

\DIFaddend \noindent\textbf{Verifiable AI.}
Matrix multiplication is also a dominant cost in verifiable machine-learning
systems. Systems such as vCNN~\cite{lee2024vcnn}, pvCNN~\cite{pvCNN}, and
zkVC~\cite{zhang2025zkvc} optimize verifiable inference for convolutional neural
networks by reducing the number of multiplication constraints induced by
QAP-based representations. Other systems use sum-check-style techniques.
zkCNN~\cite{liu2021zkcnn} proposes a sum-check-based protocol for proving
convolutions with linear prover time. In addition, zkGPT~\cite{zkGPT} and
zkLLM~\cite{zkLLM} introduce protocols specialized for LLM architectures and
use sum-check protocols to prove the correctness of linear layers efficiently.
Meanwhile, Mystique~\cite{weng2021mystique} and zkML~\cite{chen2024zkml} use
Freivalds' randomized checks. \\

\noindent\textbf{Code-based proximity testing.}
\(\protocol\) also builds on the use of linear error-correcting codes and proximity
tests in succinct proof systems. FRI~\cite{FRI} and related
works~\cite{STIR,WHIR} construct proximity IOPs for
Reed--Solomon codewords. Recent work further extends this viewpoint to more
general foldable codes~\cite{basefold}.
Our proximity layer is closer in spirit to the local codeword tests used in
Ligero~\cite{Ligero}, Brakedown~ \DIFaddbegin \DIFadd{\cite{golovnev2023brakedown}}\DIFaddend , and
Blaze~\cite{blaze}. In particular, our use of in-circuit encoding checks is
similar in spirit to Orion~\cite{orion} . \DIFaddbegin \DIFadd{Follow-up work on Orion's
soundness~\cite{restoring_orion} illustrates why codeword-validity and
proximity conditions must be discharged explicitly rather than left as input
assumptions.  Our analysis uses the proximity-gap viewpoint for
Reed--Solomon and interleaved codes~\cite{bensasson2023proximity,diamond2024interleaved}
to state the input-extraction error separately from the sampled relation error.
We do not claim novelty for code-based proximity testing itself;
our contribution is a matrix-specific reduction that combines these tools with
Freivalds' check and commitment linking to reduce in-circuit matrix-product
constraints.
}\DIFaddend 

% \noindent\textbf{Verifiable machine learning and general VC.}
% Verifiable ML systems such as SafetyNets~\cite{ghodsi2017safetynets},
% Mystique~\cite{weng2021mystique}, zkCNN~\cite{liu2021zkcnn},
% vCNN~\cite{lee2024vcnn}, and zkML~\cite{chen2024zkml} all encounter matrix
% multiplication as a major arithmetic cost.
% These systems usually verify an entire inference pipeline, either by
% arithmetizing the computation directly or by using a general proof-system
% backend for arithmetic circuits.
% \(\protocol\) targets a narrower interface: it verifies one committed matrix-product
% claim, or a batch of such claims, over committed encoded roots.
% It is therefore a checking layer that can be composed with an ML or VC system
% that already maintains committed inputs, rather than a replacement for
% end-to-end zkML or general-purpose verifiable computation.\\


% \noindent\textbf{Sum-check and GKR-based verification.}
% The sum-check protocol~\cite{lund1992algebraic} and the GKR
% framework~\cite{goldwasser2015delegating} are standard tools for verifying
% arithmetic circuits and matrix-heavy computations through multilinear
% extensions and layered circuit checks.
% Subsequent systems improve the prover/verifier trade-offs and make these
% techniques practical in SNARK-like settings~\cite{thaler2013time,
% wahby2018doubly,setty2020spartan}.
% These approaches are highly relevant to matrix multiplication, especially when
% many products are batched or embedded in a larger arithmetic circuit.
% \(\protocol\) uses a different interface: the verifier starts from committed encoded
% matrix commitments and checks local codeword equations induced by Freivalds'
% randomization.
% We therefore use Freivalds-in-SNARK as the main experimental baseline because it
% checks the same committed-root statement in the same Groth16 measurement stack.
% A direct implementation-level comparison with sum-check/GKR systems would
% require normalizing the commitment model, batching strategy, setup, transcript,
% and verifier work, so we treat it as a separate systems comparison.\\


% \noindent\textbf{Dedicated matrix-verification protocols.}
% Dedicated matrix-verification protocols, including
% zkMatrix~\cite{cong2024zkmatrix}, DualMatrix~\cite{cong2026dualmatrix},
% zkVC~\cite{zhang2025zkvc}, and coding-theoretic matrix checking by Bennett et
% al.~\cite{bennett2023matrix}, study proof systems or commitment interfaces
% specialized to matrix operations.
% These works explore trade-offs among polynomial commitments, batching, prover
% work, sparsity assumptions, and proof size.
% Our construction is closest in goal to this line, but exposes a different
% boundary: the input-certification layer is separated from the online
% multiplication checker, and the checked relation is reduced to sampled local
% fold equations over encoded commitments.
% This boundary is important for the claims in this paper: our evaluation measures the online checking layer over committed roots and does not claim to close the
%DIF >  original-matrix-to-root pipeline.


% % \section{Related Work}
% % \label{sec:related}

% % \noindent\textbf{Verifiable ML and general VC.}
% % Verifiable ML systems such as SafetyNets~\cite{ghodsi2017safetynets},
% % Mystique~\cite{weng2021mystique}, zkCNN~\cite{liu2021zkcnn},
% % vCNN~\cite{lee2024vcnn}, and zkML~\cite{chen2024zkml} all spend much of their
% % proof cost on linear layers and matrix multiplication.  Recent systems also
% % use zero-knowledge proofs for model evaluation and fairness
% % claims~\cite{south2024verifiable,zhang2025fairzk}, and CP-SNARKs have been
% % studied for zkML with committed model data~\cite{lycklama2025artemis}.  These
% % systems usually prove a full inference, training step, or model property,
% % often using a general SNARK backend such as Groth16~\cite{groth2016size}.
% % \(\protocol\) has a narrower goal: it checks one committed matrix product, or a
% % batch of such products, over committed encoded roots.  Thus, \(\protocol\) is
% % a matrix-checking layer that can be used inside a larger VC or zkML system; it
% % is not an end-to-end verifiable inference system.\\

% % \noindent\textbf{Sum-check and GKR-based verification.}
% % The sum-check protocol~\cite{lund1992algebraic} and the GKR
% % framework~\cite{goldwasser2015delegating} are standard ways to verify
% % arithmetic circuits using multilinear extensions and layered checks.  Later
% % systems improved the prover and verifier costs and made these ideas practical
% % in SNARK-like settings, including time-optimal circuit-evaluation
% % protocols~\cite{thaler2013time}, doubly-efficient
% % SNARKs~\cite{wahby2018doubly}, and Spartan~\cite{setty2020spartan}.  These
% % systems are natural when many matrix operations are part of one large
% % arithmetic circuit.  \(\protocol\) takes a different route.  It starts from
% % committed encoded matrix roots and checks local codeword equations produced by
% % Freivalds' randomization.  We therefore use Freivalds-in-SNARK as the main
% % experimental baseline: it checks the same committed-root statement with the
% % same Groth16 backend.  A direct comparison with sum-check/GKR systems would
% % also need to match the commitment model, batching method, setup, transcript,
% % and verifier work.\\

% % \noindent\textbf{Dedicated matrix-verification protocols.}
% % Classical randomized matrix checking starts with Freivalds'
% % algorithm~\cite{freivalds1977probabilistic}, which reduces a matrix-product
% % claim to a randomized vector check.  This idea also appears in SNARK-based
% % verifiable ML baselines~\cite{ghodsi2017safetynets}.  Recent dedicated
% % matrix-verification protocols include
% % zkMatrix~\cite{cong2024zkmatrix}, DualMatrix~\cite{cong2026dualmatrix},
% % zkVC~\cite{zhang2025zkvc}, and coding-theoretic matrix checking by Bennett et
% % al.~\cite{bennett2023matrix}.  These works study proof systems and commitment
% % methods for matrix operations, with different choices for batching, prover
% % work, sparsity, and proof size.  \(\protocol\) is closest to this line of
% % work.  The main difference is which costs we include.  We separate input
% % certification from the online multiplication checker, and the online checker
% % tests only sampled fold equations over encoded commitments.  Our evaluation
% % therefore measures the online checking layer over committed roots; it does not
%DIF >  % claim to measure the full cost of building the roots from original matrices.\\

% % \noindent\textbf{Technical ingredients.}
% % \(\protocol\) also uses tools from code-based proximity testing and
% % commitment linking.  These are technical ingredients, not the main comparison
% % class.  Code-based proof systems such as Ligero~\cite{ames2017ligero},
% % FRI-based protocols~\cite{ben2018fast}, Orion~\cite{xie2022orion}, and
% % Brakedown~\cite{golovnev2023brakedown} show how committed codewords and local
% % checks can support succinct verification.  Commit-and-prove
% % SNARKs~\cite{gnark,legosnark}, recent CP-SNARK systems for
% % zkML~\cite{lycklama2025artemis}, and QA-NIZK-style linking
% % proofs~\cite{jutla2013shorter,qa-nizk} provide ways to prove that sampled
% % SNARK witnesses match externally committed data.


%% Section 3: Preliminaries
\section{Preliminaries}
\label{sec:prelim}

 \DIFaddbegin \DIFadd{Throughout, \(\mathbb{F}\) is a prime field and \(\lambda\) is }\DIFaddend the security
parameter.   Matrices are denoted by bold uppercase letters  and vectors by bold
lowercase letters \DIFaddbegin \DIFadd{.  For \(\mat{M}\)}\DIFaddend , we write  \DIFaddbegin \DIFadd{\(\mat{M}[i,*]\) and
\(\mat{M}[*,j]\) for its \(i\)}\DIFaddend -th  \DIFaddbegin \DIFadd{row and \(j\)}\DIFaddend -th column.  We  \DIFaddbegin \DIFadd{use
\([n]=\{0,\ldots,n-1\}\), \(I\sample[n]^t\) for \(t\) uniform query
positions, and \(r\sample\mathbb{F}\) for a uniform field element}\DIFaddend .

\subsection{Linear  Codes \DIFaddbegin \DIFadd{and Folding}\DIFaddend }
\label{sec:prelim-codes}

 \DIFaddbegin \DIFadd{Let \(\mathcal{C}:\mathbb{F}^k\to\mathbb{F}^n\) be a linear code with
generator matrix \(\mat{G}\), rate \(\rho=k/n\), }\DIFaddend and  \DIFaddbegin \DIFadd{relative distance
\(\delta\).  We write \(\enc(\vect{m})=\vect{m}\mat{G}\).  A matrix
\(\mat{M}\in\mathbb{F}^{k\times k}\) is encoded row-wise as
\(\widehat{\mat{M}}=\enc(\mat{M})\in\mathbb{F}^{k\times n}\).
}\DIFaddend 

\begin{definition}[Interleaved code]\label{def:interleaved-code}
 \DIFaddbegin \DIFadd{The }\DIFaddend \(\ell\)-fold interleaved code  \DIFaddbegin \DIFadd{\(\mathcal{C}^{\langle\ell\rangle}\)
groups }\DIFaddend \(\ell\) codewords  position by position.   \DIFaddbegin \DIFadd{Its }\DIFaddend \(j\)-th symbol is  \DIFaddbegin \DIFadd{an
element of \(\mathbb{F}^{\ell}\).  In particular, a }\DIFaddend row-wise  \DIFaddbegin \DIFadd{encoded
\(k\times k\) matrix is a }\DIFaddend \(k\) \DIFaddbegin \DIFadd{-fold interleaved codeword whose }\DIFaddend \(j\)-th
 \DIFaddbegin \DIFadd{symbol is \(\widehat{\mat{M}}[*,j]\)}\DIFaddend .
\end{definition}

 \DIFaddbegin \DIFadd{For coefficients \(r_0,\ldots,r_{k-1}\in\mathbb{F}\), code linearity gives
}\begin{equation}\DIFadd{\label{eq:code-linearity}
  \sum_{i=0}^{k-1}r_i\widehat{\mat{M}}[i,*]
  =\enc\!\left(\sum_{i=0}^{k-1}r_i\mat{M}[i,*]\right).
}\end{equation}\DIFaddend 
 \DIFaddbegin \DIFadd{For \(\vect{r},\vect{c}\in\mathbb{F}^k\), define
}\begin{equation}\DIFadd{\label{def:fold}
  \fold(\vect{r},\vect{c})\coloneqq
  \langle\vect{r},\vect{c}\rangle.
}\end{equation}\DIFaddend 
 \DIFaddbegin \DIFadd{Thus \(\fold(\vect{r},\widehat{\mat{M}}[*,j])\) is the \(j\)}\DIFaddend -th  \DIFaddbegin \DIFadd{coordinate
of \(\enc(\vect{r}\mat{M})\).  Our implementation uses a Reed--Solomon code;
Appendix~\ref{app:rs-barycentric} gives its concrete consistency check}\DIFaddend .

\subsection{Commitment  \DIFaddbegin \DIFadd{and Proof Interfaces}\DIFaddend }
\label{sec:prelim-commitments}
 \DIFaddbegin 

\DIFadd{We use Merkle trees as position-binding vector }\DIFaddend commitments and Pedersen
 \DIFaddbegin \DIFadd{commitments as hiding, computationally binding, and additively homomorphic
commitments.  We write \(\merkleC\), \(\merkleO\), and \(\merkleV\) for
Merkle commitment, opening, and verification, }\DIFaddend and  \DIFaddbegin \DIFadd{\(\pedC\) for Pedersen
commitment}\DIFaddend .   \DIFaddbegin \DIFadd{Merkle leaves in our construction are Pedersen commitments}\DIFaddend ,  \DIFaddbegin \DIFadd{so
the tree itself uses deterministic leaves}\DIFaddend .

 \DIFaddbegin \label{sec:prelim-snark}
\DIFaddend A commit-and-prove SNARK (CP-SNARK)~ \DIFaddbegin \DIFadd{\cite{gnark,legosnark} proves an NP
relation while binding a designated witness part \(\comwit\) to a public
witness commitment \(\widetilde{\cm}=\bPi_{\cp}.\com(\para,\comwit;
\widetilde{o})\).  We use the interface
\(\bPi_{\cp}=(\setup,\com,\prove,\verify)\)}\DIFaddend .

 \label{sec:prelim-cplink}
 CP-Link  \DIFaddbegin \DIFadd{proves that a }\DIFaddend SNARK-side commitment to an ordered concatenation of
sampled blocks  \DIFaddbegin \DIFadd{contains the same messages as the corresponding external
Pedersen commitments.  For blocks \((\vect{m}_i)_{i=0}^{q-1}\), it links
}\[
\DIFadd{\begin{aligned}
  \widetilde{\cm}
  &=\com(\ck_S,\vect{m}_0\|\cdots\|\vect{m}_{q-1};\widetilde{o}),\\
  \cm_i&=\com(\ck_E,\vect{m}_i;o_i)\quad(i\in[q]).
\end{aligned}
}\]\DIFaddend 
The  \DIFaddbegin \DIFadd{construction requires }\DIFaddend CP-SNARK  \DIFaddbegin \DIFadd{soundness and zero knowledge, commitment
binding and hiding, }\DIFaddend and  CP-Link soundness  \DIFaddbegin \DIFadd{and }\DIFaddend witness zero knowledge \DIFaddbegin \DIFadd{; these
assumptions are summarized in Definition~\ref{def:lamp-backend-assumptions}.
}\DIFaddend 

\subsection{\DIFaddbegin \DIFadd{Structured }\DIFaddend Freivalds  \DIFaddbegin \DIFadd{Check}\DIFaddend }
\label{sec:prelim-freivalds}

Freivalds' algorithm~\cite{freivalds1977probabilistic}  \DIFaddbegin \DIFadd{verifies
\(\mat{A}\mat{B}=\mat{C}\) by checking
\(\vect{r}\mat{A}\mat{B}=\vect{r}\mat{C}\).  We use the compact structured
challenge \(\vect{r}=(1,r,\ldots,r^{k-1})\) for
\(r\sample\mathbb{F}\).  If
\(\mat{D}=\mat{A}\mat{B}-\mat{C}\neq0\), a nonzero column of \(\mat{D}\)
defines a nonzero }\DIFaddend univariate polynomial in  \DIFaddbegin \DIFadd{\(r\) }\DIFaddend of degree at most  \DIFaddbegin \DIFadd{\(k-1\).
By Schwartz--Zippel~\cite{schwartz1980fast,zippel1979probabilistic},
}\[
  \DIFadd{\Pr[\vect{r}\mat{D}=0]\le\frac{k-1}{|\mathbb{F}|}.
}\]
\DIFaddend 


%% Section 4: Technical Overview
\section{Overview of LAMP}
\label{sec:tech-overview}

This section gives a technical overview of \(\protocol\) before the formal
construction in Section~\ref{sec:construction}.  The main message is that
\(\protocol\) does not ask the SNARK circuit to compute a matrix product.
Instead, the prover commits to encoded data, and the circuit checks only a
small number of local consistency equations at randomly sampled codeword
positions.

\subsection{Problem Statement and Relation Reduction Chain}
\label{sec:overview-reduction-chain}

The starting point is the matrix multiplication relation
\[
  \mat{A}\mat{B}=\mat{C},
  \qquad
  \mat{A},\mat{B},\mat{C}\in\mathbb{F}^{k\times k}.
\]
A direct SNARK arithmetization of this relation computes every entry of
\(\mat{A}\mat{B}\).  This requires \(k^3\) scalar multiplications and therefore
\(\mathcal{O}(k^3)\) arithmetic constraints.  This is the original relation:
\[
\relation_{\mathsf{orig}}
=
\left\{
(\mat{A},\mat{B},\mat{C}) :
\mat{A}\mat{B}=\mat{C}
\right\}.
\]

Freivalds' algorithm reduces the cubic check to a vector-matrix check.  The
classical algorithm samples a challenge vector \(\vect{r}\in\mathbb{F}^k\) and checks
  $\vect{r}\mat{A}\mat{B}
  =
  \vect{r}\mat{C}.$

In \(\protocol\), the verifier derives one scalar \(r\) and uses the structured
vector \((1,r,\ldots,r^{k-1})\), which keeps the public challenge compact.
\DIFaddbegin \DIFadd{At the same transcript point, it derives an independent scalar \(s\) and the
structured vector \(\vect{s}=(1,s,\ldots,s^{k-1})\).  This second challenge is
used only to test the committed encoding of \(\mat{B}\).
}\DIFaddend Equivalently, the prover introduces intermediate vectors
\[
  \vect{x}=\vect{r}\mat{A},
  \qquad
  \vect{y}=\vect{x}\mat{B},
  \qquad
  \vect{z}=\vect{r}\mat{C},
\]
and the verifier checks \(\vect{y}=\vect{z}\).  This gives the Freivalds
relation
\[
\begin{aligned}
\relation_{\mathsf{Fr}}
=
\Big\{
&(\mat{A},\mat{B},\mat{C},\vect{x},\vect{y},\vect{z}) :
\\
&\vect{x}=\vect{r}\mat{A},\quad
  \vect{y}=\vect{x}\mat{B},\quad
  \vect{z}=\vect{r}\mat{C},\quad
  \vect{y}=\vect{z}
\Big\}.
\end{aligned}
\]
This removes the cubic matrix-matrix computation, but a SNARK circuit that
computes \(\vect{r}\mat{A}\), \(\vect{x}\mat{B}\), and \(\vect{r}\mat{C}\)
directly still performs \(\mathcal{O}(k^2)\) scalar multiplications.

\(\protocol\) applies one more reduction.  Rather than computing the full
vector-matrix products inside the circuit, it checks the products only at
randomly sampled positions in an encoded domain.  The prover performs the
expensive native computation outside the SNARK, commits to the encoded
matrices and intermediate vectors, and the SNARK checks local equations of the
form
\[
  \enc(\vect{x})[j]
  =
  \fold(\vect{r},\widehat{\mat{A}}[*,j]),
\]
\[
\begin{aligned}
  \enc(\vect{y})[j]
  &=
  \fold(\vect{x},\widehat{\mat{B}}[*,j]),\\
  \enc(\vect{z})[j]
  &=
  \fold(\vect{r},\widehat{\mat{C}}[*,j]).
\end{aligned}
\]
\DIFaddbegin \DIFadd{The prover additionally computes \(\vect{b}=\vect{s}\mat{B}\), encodes it as
\(\widehat{\vect{b}}=\enc(\vect{b})\), and the circuit checks
}\[
  \DIFadd{\widehat{\vect{b}}[j]
  =
  \fold(\vect{s},\widehat{\mat{B}}[*,j]).
}\]
\DIFadd{This additional fold tests the committed encoding of \(\mat{B}\) independently
of the product relation.
}\DIFaddend Only the queried columns
\(\widehat{\mat{A}}[*,j]\), \(\widehat{\mat{B}}[*,j]\),
\(\widehat{\mat{C}}[*,j]\) are used by the circuit.  Thus, for \(t\) queried
positions, the fold part of the circuit has cost \(\mathcal{O}(tk)\), which is
linear in \(k\) for fixed \(t\).

\subsection{Encoding and Proximity Testing}
\label{sec:overview-encoding}

Let \(\mathcal{C}:\mathbb{F}^k\to\mathbb{F}^n\) be a linear code
with relative distance \(\delta\).  For a matrix
\(\mat{M}\in\mathbb{F}^{k\times k}\), the prover row-wise encodes \(\mat{M}\)
as $\widehat{\mat{M}}=\enc(\mat{M})\in\mathbb{F}^{k\times n}$.
By Definition~\ref{def:interleaved-code}, this row-wise encoding can also be
viewed as a \(k\)-fold interleaved codeword, where the \(j\)-th interleaved
symbol is the encoded column \(\widehat{\mat{M}}[*,j]\).

The key identity is the linearity of the code.  For any vector
\(\vect{u}\in\mathbb{F}^k\),
\[
  \big(
  \fold(\vect{u},\widehat{\mat{M}}[*,1]),
  \ldots,
  \fold(\vect{u},\widehat{\mat{M}}[*,n])
  \big)
  =
  \enc(\vect{u}\mat{M}).
\]
Therefore, one coordinate of the encoded vector-matrix product can be checked
using one encoded column.  For example, if
\(\vect{x}=\vect{r}\mat{A}\), then for every \(j\in[n]\),
\[
  \widehat{\vect{x}}[j]
  =
  \fold(\vect{r},\widehat{\mat{A}}[*,j]).
\]
Similarly, if \(\vect{y}=\vect{x}\mat{B}\) and
\(\vect{z}=\vect{r}\mat{C}\), then
\[
  \widehat{\vect{y}}[j]
  =
  \fold(\vect{x},\widehat{\mat{B}}[*,j]),
  \qquad
  \widehat{\vect{z}}[j]
  =
  \fold(\vect{r},\widehat{\mat{C}}[*,j]).
\]

The prover commits to the encoded objects before the query positions are
sampled.  After the query set \(I=(j_1,\ldots,j_t)\) is derived, the SNARK
checks the above equations only for \(j\in I\), together with the equality
\[
  \widehat{\vect{y}}[j]=\widehat{\vect{z}}[j].
\]
The SNARK also checks that
 \DIFaddbegin \DIFadd{\(\widehat{\vect{x}},\widehat{\vect{y}},\widehat{\vect{z}},
\widehat{\vect{b}}\) }\DIFaddend are valid
encodings of the intermediate vectors
 \DIFaddbegin \DIFadd{\(\vect{x},\vect{y},\vect{z},\vect{b}\)}\DIFaddend . These code-consistency checks are separate from
the sampled fold checks and prevent the prover from using arbitrary length-\(n\)
words as the intermediate encoded views.

This is a proximity test: if a claimed vector relation is false, then the two
corresponding codewords are distinct.  Since the code has relative distance
\(\delta\), distinct codewords disagree on at least a \(\delta\)-fraction of
positions.  Random sampling therefore detects an inconsistent relation with
probability governed by \(\delta\) and the number of sampled positions.

\subsection{Sampled-Opening Layer}
\label{sec:overview-commit-link}

The sampled checks above are meaningful only if the values used inside the
SNARK are bound to data fixed before the relevant challenges are sampled.
\(\protocol\) achieves this through a sampled-opening layer.
The prover first commits to the interleaved symbols, equivalently the encoded
matrix columns, and fixes a Merkle root before the Freivalds  \DIFaddbegin \DIFadd{and independent
input-proximity challenges are }\DIFaddend derived.  After the  \DIFaddbegin \DIFadd{challenges are }\DIFaddend fixed, the prover computes the intermediate
vectors, commits to their encoded symbols, and fixes a second Merkle root before
the sampled positions are derived.

% The SNARK circuit itself only checks field equations on the sampled values.
The SNARK  \DIFaddbegin \DIFadd{checks only the sampled field equations and intermediate-vector
encoding checks.  Merkle openings authenticate }\DIFaddend the external Pedersen  \DIFaddbegin \DIFadd{leaves,
and }\DIFaddend CP-Link  \DIFaddbegin \DIFadd{binds those leaves to the sampled blocks used by the CP-SNARK.
Consequently, the expensive authentication and linking checks remain outside
the circuit.  }\DIFaddend Section~\ref{sec:security}  \DIFaddbegin \DIFadd{gives the formal soundness argument}\DIFaddend .


%% Section 5: Our Construction
\section{The LAMP Protocol}
\label{sec:construction}

Section~\ref{sec:tech-overview} reduced the original matrix-multiplication
relation to the Freivalds relation and then to sampled encoded checks.  This
section formalizes the relation checked by \(\protocol\).  The important point is
that \(\protocol\) is not a single arithmetic-circuit relation: the arithmetic
checks are proved by an inner CP-SNARK, while Merkle membership and CP-Link are
verified outside the circuit.

\subsection{The Composite LAMP Relation}
\label{sec:relation-chain}

Let \(\mat{A},\mat{B},\mat{C}\in\mathbb{F}^{k\times k}\), and let
\(\mathcal{C}:\mathbb{F}^k\to\mathbb{F}^n\) be a linear code with
relative distance \(\delta\).  The target matrix-multiplication relation is
\[
  \mathcal{R}_{\mathsf{orig}}
  =
  \{(\mat{A},\mat{B},\mat{C}) :
    \mat{A}\mat{B}=\mat{C}\}.
\]
The verifier does not receive the  \DIFaddbegin \DIFadd{original }\DIFaddend matrices as public input.  Instead, the
prover commits to row-wise encodings of the matrices and proves that the
committed objects satisfy the multiplication relation.

The relation checked by \(\protocol\) is the composite relation
\[
  \mathcal{R}_{\mathsf{LAMP}}
  =
  \mathcal{R}_{\mathsf{LAMP}}^{\mathsf{on}}
  \land
  \mathcal{R}_{\mathsf{LAMP}}^{\mathsf{off}} .
\]
Here \(\mathcal{R}_{\mathsf{LAMP}}^{\mathsf{on}}\) is the online relation proved by the CP-SNARK after the Freivalds challenge, query positions, and code-check points are sampled.  The relation \(\mathcal{R}_{\mathsf{LAMP}}^{\mathsf{off}}\) is the offline relation: it captures the values fixed by commitments before the online queries are known, and the verifier checks it through Merkle openings and CP-Link proofs.
Sections~\ref{sec:snark-relation} and~\ref{sec:openings-and-cplink} define
these two relations explicitly.
The circuit checks field equations on sampled encoded values, while the auxiliary verifier checks that those values are the ones bound by the pre-query commitments.

\DIFaddbegin \noindent\DIFadd{\textbf{Using the original matrices in an application circuit.}
The standard construction avoids placing all matrix entries in the CP-SNARK
witness: the external roots bind the encoded matrices, and the CP-SNARK commits
only to the \(\mathcal{O}(tk)\) sampled values used online.  If an application
circuit already uses \(\mat{A},\mat{B},\mat{C}\), it can instead commit to all
\(\mathcal{O}(k^2)\) entries and reuse those witness variables in the
\(\protocol\) checks.  This removes the separate Merkle/CP-Link bridge to the
application circuit while preserving
\(\mathcal{O}(tk+E_{\mathsf{code}})\) multiplication-check constraints; binding
the commitment to the intended application data remains application-specific.
}

\DIFaddend % \subsection{Commitment Layout and Public Challenges}
% \label{sec:commitment-challenge}
% The prover first encodes the matrices row-wise:
% {\footnotesize
% \(\widehat{\mat{A}}=\enc(\mat{A}),
%   \widehat{\mat{B}}=\enc(\mat{B}),
%   \widehat{\mat{C}}=\enc(\mat{C})\).
% }
% For each codeword position \(j\in[n]\), it packs the three encoded columns into
% {\footnotesize
% \(\vect{m}_{ABC,j}
%   \coloneqq
%   \big(
%     \widehat{\mat{A}}[*,j]\,
%     \|\,\widehat{\mat{B}}[*,j]\,
%     \|\,\widehat{\mat{C}}[*,j]
%   \big)
%   \in\mathbb{F}^{3k}\).
% }
% Using a Pedersen commitment key \(\ck_{ABC}\), the prover samples
% \(o_{ABC,j}\sample\mathbb{F}\) and computes
% {\footnotesize
% \(\cm_{ABC,j}
%   \gets
%   \pedC(\ck_{ABC},\vect{m}_{ABC,j};o_{ABC,j})\).
% }
% The matrix root is the Merkle commitment to these Pedersen commitments:
% {\footnotesize
% \(\rt_{ABC}
%   \gets
%   \merkleC((\cm_{ABC,j})_{j\in[n]};\bot)\).
% }
% The root \(\rt_{ABC}\) is fixed before the Freivalds challenge is derived.
% After receiving \(\rt_{ABC}\), the verifier samples \(r\sample\mathbb{F}\). 
% The prover sets {\footnotesize \(\vect{r}=(1,r,r^2,\ldots,r^{k-1})\)}
% and computes
% {\footnotesize
% \(\vect{x}=\vect{r}\mat{A},
%   \vect{y}=\vect{x}\mat{B},
%   \vect{z}=\vect{r}\mat{C}\).
% }
% It then encodes the intermediate vectors:
% {\footnotesize
% \(\widehat{\vect{x}}=\enc(\vect{x}),
%   \widehat{\vect{y}}=\enc(\vect{y}),
%   \widehat{\vect{z}}=\enc(\vect{z})\).
% }
% For each \(j\in[n]\), the prover packs the encoded vector symbols into
% {\footnotesize
% \(\vect{m}_{XYZ,j}
%   \coloneqq
%   \big(
%     \widehat{\vect{x}}[j]\,
%     \|\,\widehat{\vect{y}}[j]\,
%     \|\,\widehat{\vect{z}}[j]
%   \big)
%   \in\mathbb{F}^{3}\).
% }
% Using a Pedersen commitment key \(\ck_{XYZ}\), the prover samples
% \(o_{XYZ,j}\sample\mathbb{F}\), computes
% {\footnotesize
% \(\cm_{XYZ,j}
%   \gets
%   \pedC(\ck_{XYZ},\vect{m}_{XYZ,j};o_{XYZ,j})\),
% }
% and sends the vector-symbol root
% {\footnotesize
% \(\rt_{XYZ}
%   \gets
%   \merkleC((\cm_{XYZ,j})_{j\in[n]};\bot)\)
% }
% to the verifier.
% Only after receiving \(\rt_{XYZ}\) does the verifier sample the query tuple
% {\footnotesize \(I=(j_1,\ldots,j_t)\sample [n]^t\).}
% It also samples public code-check points
% {\footnotesize
% \(\alpha_x,\alpha_y,\alpha_z\sample\mathbb{F}\setminus D\),
% }
% where \(D\) is the
% code evaluation domain.  These points are used by the circuit to check that
% {\footnotesize
% \(\widehat{\vect{x}},\widehat{\vect{y}},\widehat{\vect{z}}\)
% }
% are encodings of
% the claimed intermediate vectors.  For the Reed--Solomon code used in our implementation, this check is the barycentric out-of-domain consistency test recalled in Appendix~\ref{app:rs-barycentric}. 
\subsection{Commitment Layout and Public Challenges}
\label{sec:commitment-challenge}

The prover first encodes the matrices row-wise:
{\footnotesize
\(\widehat{\mat{A}}=\enc(\mat{A}),
  \widehat{\mat{B}}=\enc(\mat{B}),
  \widehat{\mat{C}}=\enc(\mat{C})\).
}
For each codeword position \(j\in[n]\), it packs the three \(k\)-dimensional
interleaved symbols, equivalently the three encoded columns, into
{\footnotesize
\(\vect{m}_{ABC,j}
  \coloneqq
  \big(
    \widehat{\mat{A}}[*,j]\,
    \|\,\widehat{\mat{B}}[*,j]\,
    \|\,\widehat{\mat{C}}[*,j]
  \big)
  \in\mathbb{F}^{3k}\).
}
Using a Pedersen commitment key \(\ck_{ABC}\), the prover samples
\(o_{ABC,j}\sample\mathbb{F}\) and computes
{\footnotesize
\(\cm_{ABC,j}
  \gets
  \pedC(\ck_{ABC},\vect{m}_{ABC,j};o_{ABC,j})\).
}
The matrix root is the Merkle commitment to these Pedersen commitments:
{\footnotesize
\(\rt_{ABC}
  \gets
  \merkleC((\cm_{ABC,j})_{j\in[n]};\bot)\).
}
Here and below, \(\bot\) indicates that no additional randomness is used in the
Merkle tree, since the leaves are Pedersen commitments that already provide
hiding. The root \(\rt_{ABC}\) is fixed before the Freivalds challenge is
derived.

After receiving \(\rt_{ABC}\), the verifier  \DIFaddbegin \DIFadd{independently samples
\(r,s\sample\mathbb{F}\)}\DIFaddend .  The prover sets
{\footnotesize
\(\vect{r}=(1,r,r^2,\ldots,r^{k-1})\) \DIFaddbegin \DIFadd{and
\(\vect{s}=(1,s,s^2,\ldots,s^{k-1})\)}\DIFaddend }
and computes
{\footnotesize
 \DIFaddbegin \DIFadd{\(\vect{x}=\vect{r}\mat{A},
  \vect{y}=\vect{x}\mat{B},
  \vect{z}=\vect{r}\mat{C},
  \vect{b}=\vect{s}\mat{B}\)}\DIFaddend .
}
It then encodes the intermediate vectors:
{\footnotesize
 \DIFaddbegin \DIFadd{\(\widehat{\vect{x}}=\enc(\vect{x}),
  \widehat{\vect{y}}=\enc(\vect{y}),
  \widehat{\vect{z}}=\enc(\vect{z}),
  \widehat{\vect{b}}=\enc(\vect{b})\)}\DIFaddend .
}
The encoded intermediate vectors form a  \DIFaddbegin \DIFadd{\(4\)}\DIFaddend -fold interleaved codeword.  For
each \(j\in[n]\), the prover packs its \(j\)-th interleaved symbol into
{\footnotesize
 \DIFaddbegin \DIFadd{\(\vect{m}_{XYZB,j}
  \coloneqq
  \big(
    \widehat{\vect{x}}[j]\,
    \|\,\widehat{\vect{y}}[j]\,
    \|\,\widehat{\vect{z}}[j]\,
    \|\,\widehat{\vect{b}}[j]
  \big)
  \in\mathbb{F}^{4}\)}\DIFaddend .
}
Using a Pedersen commitment key  \DIFaddbegin \DIFadd{\(\ck_{XYZB}\)}\DIFaddend , the prover samples
 \DIFaddbegin \DIFadd{\(o_{XYZB,j}\sample\mathbb{F}\)}\DIFaddend , computes
{\footnotesize
 \DIFaddbegin \DIFadd{\(\cm_{XYZB,j}
  \gets
  \pedC(\ck_{XYZB},\vect{m}_{XYZB,j};o_{XYZB,j})\)}\DIFaddend ,
}
and sends the vector-symbol root
{\footnotesize
 \DIFaddbegin \DIFadd{\(\rt_{XYZB}
  \gets
  \merkleC((\cm_{XYZB,j})_{j\in[n]};\bot)\)
}\DIFaddend }
to the verifier.

\DIFaddbegin \DIFadd{The fourth vector supplies an independent fold of the committed encoding of
\(\mat{B}\); its soundness role is analyzed in Section~\ref{sec:security}.
}

\DIFaddend Only after receiving  \DIFaddbegin \DIFadd{\(\rt_{XYZB}\) }\DIFaddend does the verifier sample the query tuple
{\footnotesize \(I=(j_1,\ldots,j_t)\sample [n]^t\).}
It also samples public code-check points
{\footnotesize
 \DIFaddbegin \DIFadd{\(\alpha_x,\alpha_y,\alpha_z,\alpha_b\sample\mathbb{F}\setminus D\)}\DIFaddend ,
}
where \(D\) is the
code evaluation domain. These points are used by the circuit to check that
{\footnotesize
 \DIFaddbegin \DIFadd{\(\widehat{\vect{x}},\widehat{\vect{y}},\widehat{\vect{z}},
\widehat{\vect{b}}\)
}\DIFaddend }
are encodings of
the claimed intermediate vectors. For the Reed--Solomon code used in our
implementation, this check is the barycentric out-of-domain consistency test
recalled in Appendix~\ref{app:rs-barycentric}.


\subsection{The Online Relation}
\label{sec:snark-relation}
The online relation is the arithmetic relation proved inside the
inner CP-SNARK.  Its role is not to verify Merkle paths or Pedersen openings.
Instead, it checks only the field equations that certify the sampled encoded
positions are consistent with the Freivalds reduction.

The public statement contains the Merkle roots and verifier-sampled
challenges:
 \DIFaddbegin \[
\DIFadd{\stmt = \big( \rt_{ABC},\rt_{XYZB},r,s,I,\alpha_x,\alpha_y,\alpha_z,\alpha_b
  \big).
}\]\DIFaddend 
The roots are included in the statement to bind the CP-SNARK proof to the same
transcript and offline checks; the online field equations themselves do not
evaluate Merkle paths.

The complete SNARK witness is
  $\wit=(\comwit,\omega)$,
where the committed part of the witness is
 \DIFaddbegin \[
\DIFadd{\comwit = \big( (\vect{m}_{ABC,j})_{j\in I},(\vect{m}_{XYZB,j})_{j\in I} \big),
}\]\DIFaddend 
and the remaining private witness is
   \DIFaddbegin \DIFadd{$\omega
  =
  \big(
    \vect{x},\vect{y},\vect{z},\vect{b},
    \widehat{\vect{x}},\widehat{\vect{y}},\widehat{\vect{z}},
    \widehat{\vect{b}}
  \big).$
}\DIFaddend The CP-SNARK treats \(\comwit\) as the committed part of the witness.  The
prover first computes the SNARK-side witness commitments
\(\widetilde{\cm}_{ABC}\) and  \DIFaddbegin \DIFadd{\(\widetilde{\cm}_{XYZB}\) }\DIFaddend using
\(\bPi_{\cp}.\com\), and then proves the online relation with witness
\(\wit=(\comwit,\omega)\).  These witness commitments are not checked against
the Merkle roots inside the circuit.  Instead, they are connected to the
externally committed Merkle leaves by the offline relation in
Section~\ref{sec:openings-and-cplink}.
For the CP-Link statements below, let \(\ck_{\cp,T}\) denote the public
commitment key used by the CP-SNARK witness commitment
\(\widetilde{\cm}_{T}\), for  \DIFaddbegin \DIFadd{\(T\in\{ABC,XYZB\}\)}\DIFaddend .
Write
  $\mathsf{EncCheck}_{\mathcal{C}}(\widehat{\vect{v}},\vect{v};\alpha)=1$
when the out-of-domain code-consistency check accepts that
\(\widehat{\vect{v}}\) is consistent with \(\enc(\vect{v})\) at the public
point \(\alpha\).  The online relation is defined in
Figure~\ref{relation:online}.

\begin{figure}[t]
\begin{tcolorbox}[colback=white, colframe=black]
 \DIFaddbeginFL \pseudocodeblock[mode=text]{
$\circ$ $\relation_{\mathsf{LAMP}}^{\mathsf{on}}(\stmt; \wit):$ \\
    $\qquad$parse \\
    $\qquad \qquad \stmt =
      \big(
        \rt_{ABC},\rt_{XYZB},r,s,I,\alpha_x,\alpha_y,\alpha_z,\alpha_b
      \big)$ \\
    $\qquad \qquad \wit=(\comwit,\omega)$ \\
    $\qquad \qquad \comwit =
      \big(
        (\vect{m}_{ABC,j})_{j\in I},
        (\vect{m}_{XYZB,j})_{j\in I}
      \big)$ \\
    $\qquad \qquad \omega =
      \big(
        \vect{x},\vect{y},\vect{z},\vect{b},
        \widehat{\vect{x}},\widehat{\vect{y}},\widehat{\vect{z}},
        \widehat{\vect{b}}
      \big)$ \\[0.5ex]
    $\qquad \vect{r}\gets(1,r,\ldots,r^{k-1})$ \\[0.5ex]
    $\qquad \vect{s}\gets(1,s,\ldots,s^{k-1})$ \\[0.5ex]
    $\qquad \mathsf{EncCheck}_{\mathcal{C}}
      (\widehat{\vect{x}},\vect{x};\alpha_x)=1$ \\
    $\qquad \mathsf{EncCheck}_{\mathcal{C}}
      (\widehat{\vect{y}},\vect{y};\alpha_y)=1$ \\
    $\qquad \mathsf{EncCheck}_{\mathcal{C}}
      (\widehat{\vect{z}},\vect{z};\alpha_z)=1$ \\
    $\qquad \mathsf{EncCheck}_{\mathcal{C}}
      (\widehat{\vect{b}},\vect{b};\alpha_b)=1$ \\[0.5ex]
    $\qquad \forall j\in I:$ \\
    $\qquad \qquad$parse
      $\vect{m}_{ABC,j}$ as
      $\big(
        \widehat{\mat{A}}[*,j]\,
        \|\,\widehat{\mat{B}}[*,j]\,
        \|\,\widehat{\mat{C}}[*,j]
      \big)$ \\
    $\qquad \qquad$parse
      $\vect{m}_{XYZB,j}$ as
      $\big(
        \widehat{\vect{x}}[j]\,
        \|\,\widehat{\vect{y}}[j]\,
        \|\,\widehat{\vect{z}}[j]\,
        \|\,\widehat{\vect{b}}[j]
      \big)$ \\[0.5ex]
    $\qquad \qquad
      \widehat{\vect{x}}[j]
      =
      \fold(\vect{r},\widehat{\mat{A}}[*,j])$ \\
    $\qquad \qquad
      \widehat{\vect{y}}[j]
      =
      \fold(\vect{x},\widehat{\mat{B}}[*,j])$ \\
    $\qquad \qquad
      \widehat{\vect{z}}[j]
      =
      \fold(\vect{r},\widehat{\mat{C}}[*,j])$ \\
    $\qquad \qquad
      \widehat{\vect{b}}[j]
      =
      \fold(\vect{s},\widehat{\mat{B}}[*,j])$ \\
    $\qquad \qquad
      \widehat{\vect{y}}[j]
      =
      \widehat{\vect{z}}[j]$
}
\DIFaddendFL \end{tcolorbox}
\caption{The Online Relation}
\label{relation:online}
\end{figure}
\FloatBarrier


The  \DIFaddbegin \DIFadd{encoding checks validate the four intermediate views, and the remaining
equations check the sampled folds and \(\widehat{\vect{y}}[j]=
\widehat{\vect{z}}[j]\).  Binding }\DIFaddend these sampled values  \DIFaddbegin \DIFadd{to }\DIFaddend the pre-query  \DIFaddbegin \DIFadd{roots
is the role of the offline relation}\DIFaddend .

\DIFaddbegin \subsection{\DIFadd{The Offline Relation}}
\label{sec:openings-and-cplink}

\DIFaddend For compactness, Figure~\ref{relation:lamp-off} writes the CP-Link statement
for  \DIFaddbegin \DIFadd{\(T\in\{ABC,XYZB\}\) }\DIFaddend as
\[
  \stmt_{\mathsf{link}}^T
  \coloneqq
  \big(
    (\ck_{\cp,T},\widetilde{\cm}_{T}),
    (\ck_T,\cm_{T,j})_{j\in I}
  \big),
\]
where \(\ck_{\cp,T}\) is the witness-commitment key used by the CP-SNARK
commitment \(\widetilde{\cm}_{T}\), and \(\ck_T\) is the external Pedersen
commitment key for the Merkle leaves of type \(T\).

 The offline relation connects the sampled witnesses used by the
CP-SNARK to the commitments that were fixed before the query positions were
sampled.  The term ``offline'' indicates that these checks are performed outside the SNARK circuit. In our system, it also reflects that the committed data are fixed before the online sampled checks are carried out.

For each queried position \(j\in I\), the verifier receives Merkle openings for the external Pedersen leaves $\cm_{ABC,j}$ and  \DIFaddbegin \DIFadd{$\cm_{XYZB,j}$}\DIFaddend .
The Merkle openings certify that these leaves are contained in the roots
\(\rt_{ABC}\) and  \DIFaddbegin \DIFadd{\(\rt_{XYZB}\)}\DIFaddend , respectively.  In addition, the verifier checks
CP-Link proofs connecting the CP-SNARK witness commitments
\(\widetilde{\cm}_{ABC}\) and  \DIFaddbegin \DIFadd{\(\widetilde{\cm}_{XYZB}\) }\DIFaddend to the corresponding
external Pedersen leaves.

 \DIFaddbegin \DIFadd{For each \(T\in\{ABC,XYZB\}\), }\DIFaddend CP-Link  \DIFaddbegin \DIFadd{proves that }\DIFaddend the  \DIFaddbegin \DIFadd{SNARK-side commitment
contains the ordered concatenation
\(\vect{m}_{T,j_1}\|\cdots\|\vect{m}_{T,j_t}\) committed }\DIFaddend by the corresponding
external  \DIFaddbegin \DIFadd{leaves.  Appendix~\ref{app:cplink-instant} gives the concrete
blockwise relation.
}\DIFaddend 


The offline relation is defined in Figure~\ref{relation:lamp-off}.

\begin{figure}[!t]
\begin{tcolorbox}[colback=white, colframe=black]
 \DIFaddbeginFL \pseudocodeblock[mode=text]{
$\circ\;
\relation_{\mathsf{LAMP}}^{\mathsf{off}}\big($
    $
      \rt_{ABC},\rt_{XYZB},I,
      \widetilde{\cm}_{ABC},\widetilde{\cm}_{XYZB},
      $ \\
    $\qquad
      \pi_{\mathsf{link}}^{ABC},\pi_{\mathsf{link}}^{XYZB},$ \\
    $\qquad
      (\cm_{ABC,j},\pi_{\mathsf{mem},j}^{ABC})_{j\in I},
      (\cm_{XYZB,j},\pi_{\mathsf{mem},j}^{XYZB})_{j\in I}
    \big):$ \\[0.5ex]
    $\qquad \forall j\in I:$ \\
    $\qquad \qquad
        \merkleV(
            \rt_{ABC},
            j,
            \cm_{ABC,j},
            \pi_{\mathsf{mem},j}^{ABC}
        ) = 1$ \\
    $\qquad \qquad
        \merkleV(
            \rt_{XYZB},
            j,
            \cm_{XYZB,j},
            \pi_{\mathsf{mem},j}^{XYZB}
        ) = 1$ \\[0.5ex]
    $\qquad
        \bPi_{\mathsf{link}}.\verify(
            \para_{\mathsf{link}},
            \stmt_{\mathsf{link}}^{ABC},
            \pi_{\mathsf{link}}^{ABC}
        ) = 1$ \\
    $\qquad
        \bPi_{\mathsf{link}}.\verify(
            \para_{\mathsf{link}},
            \stmt_{\mathsf{link}}^{XYZB},
            \pi_{\mathsf{link}}^{XYZB}
        ) = 1$
}
\DIFaddendFL \end{tcolorbox}
\caption{The Offline Relation}
\label{relation:lamp-off}
\end{figure}
\FloatBarrier

 \DIFaddbegin \DIFadd{Together, CP-Link and Merkle membership establish
\(\widetilde{\cm}_{T}\leftrightarrow(\cm_{T,j})_{j\in I}
\leftrightarrow\rt_T\) for \(T\in\{ABC,XYZB\}\).  Hence the }\DIFaddend CP-SNARK  \DIFaddbegin \DIFadd{checks
values bound by the pre-query roots without verifying Pedersen openings or
Merkle paths inside the circuit}\DIFaddend .

% \begin{table}[t]
% \centering
% \footnotesize
% \setlength{\tabcolsep}{4pt}
% \begin{tabularx}{\columnwidth}{@{}lcc@{}}
% \toprule
% \textbf{Component} & \textbf{Prover} & \textbf{Verifier} \\
% \midrule
% ABC encoding
% & \(E_{\enc}(k,n)\)
% & -- \\

% ABC commit
% & \(\mathcal{O}(3nk)\,\mathbb{G}+\mathcal{O}(n)\,\mathbb{H}\)
% & -- \\

% Intermediate
% & \(\mathcal{O}(k^2)\,\mathbb{F}\)
% & -- \\

% XYZ encoding
% & \(E_{\enc}(k,n)\)
% & -- \\

% XYZ commit
% & \(\mathcal{O}(n)\,\mathbb{G}+\mathcal{O}(n)\,\mathbb{H}\)
% & -- 
% \\
% Merkle openings
% & \(\mathcal{O}(t\log n) \mathbb{H}\)
% & \(\mathcal{O}(t\log n)\mathbb{H} \) 
% \\
% % CP-SNARK
% % & \(
% % \mathcal{O}(tk)+E_{\mathsf{code}} +
% % E_{\mathsf{snark}}^\mathcal{P}
% % \)
% % & \(E_{\mathsf{snark}}^\mathcal{V}\) 
% CP-SNARK
% & \(\begin{array}{c}
% \mathcal{O}(tk)+E_{\mathsf{code}}\ \text{constraints}\\
% E_{\mathsf{snark}}^\mathcal{P}
% \end{array}\)
% & \(E_{\mathsf{snark}}^\mathcal{V}\)
% \\
% CP-Link
% & \(E_{\mathsf{link}}^\mathcal{P}(k,t)\)
% & \(E_{\mathsf{link}}^\mathcal{V}(k,t)\) \\
% \bottomrule
% \end{tabularx}
% \vspace{6pt}
% \caption{Asymptotic costs for \(\protocol\).}
% \label{tab:cost-summary}
% \end{table}

\subsection{Protocol Algorithms and Complexity}
\label{sec:protocol-complexity}
Figure~\ref{fig:protocol} presents the public-coin version of
\(\protocol\).  The protocol can be made non-interactive via the
Fiat--Shamir transform, by deriving  \DIFaddbegin \DIFadd{\((r,s)\) }\DIFaddend after \(\rt_{ABC}\) and deriving
 \DIFaddbegin \DIFadd{\(I,\alpha_x,\alpha_y,\alpha_z,\alpha_b\) after \(\rt_{XYZB}\)}\DIFaddend .

The proof consists of the CP-SNARK proof, the sampled leaf commitments and
Merkle openings, and the CP-Link proofs.  The sampled matrix-column blocks and
vector-symbol blocks are private CP-SNARK witnesses; they are not serialized as
public proof material.
The public statement contains the roots
 \DIFaddbegin \DIFadd{\(\rt_{ABC},\rt_{XYZB}\) }\DIFaddend and the verifier-sampled challenges. \\

For  \DIFaddbegin \DIFadd{each \(T\in\mathcal{T}=\{ABC,XYZB\}\), the }\DIFaddend CP-Link statement  \DIFaddbegin \DIFadd{contains
the SNARK-side witness commitment and the \(t\) corresponding external leaf
commitments; the witness contains their messages and opening randomness}\DIFaddend .\\

% \paragraph{Cost model}
\noindent\textbf{Cost model.}
We report the cost of  \protocol.  The computation of the claimed output matrix \(C=AB\) is outside this cost model.
We do count the work performed by the \protocol prover after the matrices are available, including encoding, commitment generation, computation of the Freivalds intermediate witnesses, sampled Merkle openings, CP-SNARK proving, and CP-Link proving.

Let \(E_{\enc}(k,n)\) denote \DIFaddbegin \DIFadd{the }\DIFaddend cost of the corresponding encoding phase.
Let \(E_{\mathsf{code}}=E_{\mathsf{code}}(k,n)\) denote the constraint cost of the intermediate
code-consistency checks; for the Reed--Solomon check in
Appendix~\ref{app:rs-barycentric}, this is \(\mathcal{O}(k+n)\) after
domain-dependent coefficients are precomputed.
We write \(\mathbb{G}\) for group operations used by Pedersen commitments,
\(\mathbb{F}\) for field operations, and \(\mathbb{H}\) for hash evaluations.
Let \(E_{\mathsf{snark}}^{\mathcal{P}}\) and \(E_{\mathsf{snark}}^{\mathcal{V}}\) denote the CP-SNARK proving and
verification costs, respectively. Let \(E_{\mathsf{link}}^\mathcal{P}(k,t)\) and \(E_{\mathsf{link}}^\mathcal{V}(k,t)\) denote the prover and verifier costs, respectively, of the selected CP-Link construction.
For the QALink instantiation in Appendix~\ref{app:cplink-instant}, each of the
two link statements has \(q=t\) external commitments, so
\(E_{\mathsf{link}}^\mathcal{V}(k,t)\) includes \(2(t+2)\) Miller-loop inputs,
batched into one final exponentiation, plus \(\mathcal{O}(t)\) group scalar
multiplications for batching.

Table~\ref{tab:cost-summary} summarizes the asymptotic costs at the level of the main protocol components. \\

\begin{table}[ht]
\centering
\footnotesize
\setlength{\tabcolsep}{4pt}
\begin{tabularx}{\columnwidth}{@{}lcc@{}}
\toprule
\textbf{Component} & \textbf{Prover} & \textbf{Verifier} \\
\midrule
ABC encoding
& \(E_{\enc}(k,n)\)
& -- \\

ABC commit
& \(\mathcal{O}(3nk)\,\mathbb{G}+\mathcal{O}(n)\,\mathbb{H}\)
& -- \\

Intermediate
& \(\mathcal{O}(k^2)\,\mathbb{F}\)
& -- \\

 \DIFaddbeginFL \DIFaddFL{XYZB }\DIFaddendFL encoding
& \(E_{\enc}(k,n)\)
& -- \\

 \DIFaddbeginFL \DIFaddFL{XYZB }\DIFaddendFL commit
& \(\mathcal{O}(n)\,\mathbb{G}+\mathcal{O}(n)\,\mathbb{H}\)
& -- 
\\
Merkle openings
& \(\mathcal{O}(t\log n) \mathbb{H}\)
& \(\mathcal{O}(t\log n)\mathbb{H} \) 
\\
% CP-SNARK
% & \(
% \mathcal{O}(tk)+E_{\mathsf{code}} +
% E_{\mathsf{snark}}^\mathcal{P}
% \)
% & \(E_{\mathsf{snark}}^\mathcal{V}\) 
CP-SNARK
& \(\begin{array}{c}
\mathcal{O}(tk)+E_{\mathsf{code}}\ \text{constraints}\\
E_{\mathsf{snark}}^\mathcal{P}
\end{array}\)
& \(E_{\mathsf{snark}}^\mathcal{V}\)
\\
CP-Link
& \(E_{\mathsf{link}}^\mathcal{P}(k,t)\)
& \(E_{\mathsf{link}}^\mathcal{V}(k,t)\) \\
\bottomrule
\end{tabularx}
\vspace{6pt}
\caption{Asymptotic costs for \(\protocol\).}
\label{tab:cost-summary}
\end{table}


\noindent\textbf{Prover work.}
The prover first encodes the three input matrices and commits to the packed
encoded matrix columns.  The ABC commitment phase consists of \(n\) Pedersen
commitments to blocks of length \(3k\), followed by a Merkle root over
the resulting leaves.  After the  \DIFaddbegin \DIFadd{independent challenges \((r,s)\) are }\DIFaddend sampled, the
prover computes the intermediate witnesses
   \DIFaddbegin \DIFadd{$\vect{x}=\vect{r}\mat{A},
  \vect{y}=\vect{x}\mat{B},
  \vect{z}=\vect{r}\mat{C},
  \vect{b}=\vect{s}\mat{B},$
}\DIFaddend which costs \(\mathcal{O}(k^2)\) field operations.  This is distinct from the
excluded native computation of \(C=AB\).  The prover then encodes
 \DIFaddbegin \DIFadd{\(\vect{x},\vect{y},\vect{z},\vect{b}\)}\DIFaddend , commits to their encoded symbols, opens the
sampled Merkle paths, proves the online CP-SNARK relation, and produces the
CP-Link proofs.

For each queried position \(j\in I\), the CP-SNARK circuit checks  \DIFaddbegin \DIFadd{four
}\DIFaddend length-\(k\) fold equations:
 \DIFaddbegin \[
\DIFadd{\begin{aligned}
  &\fold(\vect{r},\widehat{\mat{A}}[*,j]),\quad
  \fold(\vect{x},\widehat{\mat{B}}[*,j]),\\
  &\fold(\vect{r},\widehat{\mat{C}}[*,j]),\quad
  \fold(\vect{s},\widehat{\mat{B}}[*,j]).
\end{aligned}
}\]\DIFaddend 
Thus the sampled arithmetic checks contribute \(\mathcal{O}(tk)\) constraints. 

The intermediate-vector code-consistency checks add
\(E_{\mathsf{code}}\) constraints. For the fixed-rate
Reed--Solomon setting used in the experiments, \(n=2k\), so this extra cost is
linear in \(k\) and the total online relation remains linear in \(k\) for fixed
\(t\). \\

% \paragraph{Verifier work}
\noindent\textbf{Verifier work.}
The verifier does not encode matrices, compute intermediate vectors, or build
commitment roots.  Its work consists of CP-SNARK verification, verification of
the sampled Merkle authentication paths, and CP-Link verification.  Therefore,
% using the notation above, the verifier cost is $E_{\mathsf{snark}}^\mathcal{V} + \mathcal{O}(t\log n)\mathbb{H} + E_{\mathsf{link}}^\mathcal{V}(k,t).$ \\
using the notation above, the verifier cost is
$E_{\mathsf{snark}}^\mathcal{V} + \mathcal{O}(t\log n)\mathbb{H} +
E_{\mathsf{link}}^\mathcal{V}(k,t)$, where the QALink verifier contributes
\(\mathcal{O}(t)\) pairing terms. \\

% \paragraph{Comparison}
\noindent\textbf{Comparison.}
A naive SNARK arithmetization of matrix multiplication costs
\(\mathcal{O}(k^3)\) constraints, while a Freivalds-based SNARK baseline computes
the vector-matrix products inside the circuit and costs \(\mathcal{O}(k^2)\) constraints.  
% \(\protocol\) moves the vector-matrix products to native prover work and checks only sampled encoded-column relations inside the CP-SNARK.  
% As a result, the main circuit constraint count is \(\mathcal{O}(tk)\).
\(\protocol\) moves the vector-matrix products to native prover work and checks
sampled encoded-column relations plus intermediate code-consistency equations
inside the CP-SNARK.
As a result, the main circuit constraint count is
\(\mathcal{O}(tk+E_{\mathsf{code}})\), which is linear in \(k\) for the
fixed-rate code used in our implementation.


\begin{figure*}[p]
\centering
\scriptsize
\begingroup
\setlength{\fboxsep}{2pt}
 \DIFaddbeginFL \resizebox{0.98\textwidth}{!}{%
\fbox{%
\procedure[]{\protocol}{%
    \mathcal{P}(\mat{A},\mat{B},\mat{C}) \< \< \mathcal{V} \\
    \widehat{\mat{A}},\widehat{\mat{B}},\widehat{\mat{C}} \gets \enc(\mat{A}),\enc(\mat{B}),\enc(\mat{C}) \in \mathbb{F}^{k \times n} \< \< \\
    \text{for each } j\in[n]: \\ 
    \quad \vect{m}_{ABC,j} \coloneqq \big( \widehat{\mat{A}}[*,j]\, \|\,\widehat{\mat{B}}[*,j]\, \|\,\widehat{\mat{C}}[*,j] \big) \in\mathbb{F}^{3k} \\
    \quad o_{ABC,j}\sample\mathbb{F} \\
    \quad \cm_{ABC,j} \gets \pedC(\ck_{ABC},\vect{m}_{ABC,j};o_{ABC,j}) 
    \\
    \rt_{ABC} \gets \merkleC((\cm_{ABC,j})_{j\in[n]};\bot) 
    \< \sendmessageright*[1.3cm]{\rt_{ABC}} 
    \\
    \< \sendmessageleft*[1.3cm]{r,s} \<  r,s \sample \FF 
    \\
    \vect{r}\coloneqq(1,r,\ldots,r^{k-1}) \\ 
    \vect{s}\coloneqq(1,s,\ldots,s^{k-1}) \\
    \vect{x}\gets\vect{r}\mat{A}, \quad \vect{y}\gets\vect{x}\mat{B}, \quad \vect{z}\gets\vect{r}\mat{C}, \quad \vect{b}\gets\vect{s}\mat{B}
    \\ \widehat{\vect{x}},\widehat{\vect{y}},\widehat{\vect{z}},\widehat{\vect{b}} \gets \enc(\vect{x}),\enc(\vect{y}),\enc(\vect{z}),\enc(\vect{b}) \\ 
    \text{for each } j\in[n]: \\ 
    \quad \vect{m}_{XYZB,j} \coloneqq \big( \widehat{\vect{x}}[j]\, \|\,\widehat{\vect{y}}[j]\, \|\,\widehat{\vect{z}}[j]\, \|\,\widehat{\vect{b}}[j] \big) \in\mathbb{F}^{4} \\ 
    \quad o_{XYZB,j}\sample\mathbb{F} \\ 
    \quad \cm_{XYZB,j} \gets \pedC(\ck_{XYZB},\vect{m}_{XYZB,j};o_{XYZB,j}) \\ 
    \rt_{XYZB} \gets \merkleC((\cm_{XYZB,j})_{j\in[n]};\bot) \< \sendmessageright*[1.3cm]{\rt_{XYZB}} \< \\
    \< \< I=(j_1,\ldots,j_t)\sample[n]^t 
    \\ 
    \< \sendmessageleft*[1.3cm]{I, \alpha_x, \alpha_y, \alpha_z, \alpha_b} \< \alpha_x,\alpha_y,\alpha_z,\alpha_b\sample\mathbb{F}\setminus D 
    \\
    \mathcal{T}\coloneqq(ABC,XYZB) \\
    \text{For every } T\in\mathcal{T}\text{ and every } j\in I,\\ 
    \quad \pi_{\mathsf{mem},j}^{T} \gets \merkleO((\cm_{T,\ell})_{\ell\in[n]},j) \\[0.3ex]
    \stmt \coloneqq
      \big(\rt_{ABC},\rt_{XYZB},r,s,I,\alpha_x,\alpha_y,\alpha_z,\alpha_b\big) \\
    \comwit \coloneqq
      \big((\vect{m}_{T,j})_{j\in I}\big)_{T\in\mathcal{T}}, \quad
    \omega \coloneqq
      \big(
      \vect{x},\vect{y},\vect{z},\vect{b},
      \widehat{\vect{x}},\widehat{\vect{y}},\widehat{\vect{z}},\widehat{\vect{b}}
      \big) \\
    \wit \coloneqq \big(\comwit,\omega\big) \\
    \text{For every } T\in\mathcal{T},\\
    \quad \widetilde{\cm}_{T} \gets
      \bPi_{\cp}.\com(\para,(\vect{m}_{T,j})_{j\in I};\widetilde{o}_{T}) \\
    \stmt_{\cp} \gets
      \big(
      \stmt,
      (\widetilde{\cm}_{T})_{T\in\mathcal{T}}
      \big) \\
    \pi_{\cp} \gets \bPi_{\cp}.\prove(\para,\stmt_{\cp},\wit) \\
    \text{For every } T\in\mathcal{T},\\
    \quad \pi_{\mathsf{link}}^{T} \gets
      \bPi_{\mathsf{link}}.\prove(
      \para_{\mathsf{link}},
      \stmt_{\mathsf{link}}^T,
      \wit_{\mathsf{link}}^T) \\
    \pi_{\protocol}
        \coloneqq
        \big(
        \pi_{\cp},
        (\widetilde{\cm}_{T})_{T\in\mathcal{T}},
        \{(\cm_{T,j},\pi_{\mathsf{mem},j}^{T})\}_{T\in\mathcal{T},j\in I},
(\pi_{\mathsf{link}}^{T})_{T\in\mathcal{T}}
        \big)  \< \sendmessageright*[1.3cm]{\pi_\protocol}  \<
   \text{Parse } \pi_{\protocol} 
   \\
    \< \< \stmt_{\cp} \gets \big( \stmt,(\widetilde{\cm}_{T})_{T\in\mathcal{T}} \big) \\
    \< \< b_{\cp} \gets \bPi_{\cp}.\verify \big( \para, \stmt_{\cp}, \pi_{\cp}) \\
    \< \< \text{For every } T\in\mathcal{T}: \\
    \< \< \qquad \text{For every } j\in I: \\
    \< \< \qquad \qquad b_{\mathsf{mem},j}^{T} \gets
      \merkleV( \rt_{T}, j, \cm_{T,j}, \pi_{\mathsf{mem},j}^{T}) \\
    \< \< \qquad b_{\mathsf{link}}^{T} \gets
      \bPi_{\mathsf{link}}.\verify \big(
      \para_{\mathsf{link}},
      \stmt_{\mathsf{link}}^T,
      \pi_{\mathsf{link}}^{T}
      \big) \\
    \< \< \text{Accept iff }
      b_{\cp}
      \land
      \bigwedge_{T\in\mathcal{T}}
      \left(
        b_{\mathsf{link}}^{T}
        \land
        \bigwedge_{j\in I} b_{\mathsf{mem},j}^{T}
      \right)
}
}
}
\DIFaddendFL \endgroup
\caption{The \protocol\ protocol}
\label{fig:protocol}
\end{figure*}


%% Section 6: Security Analysis
\section{Security Analysis}
\label{sec:security}

We  \DIFaddbegin \DIFadd{analyze }\DIFaddend the public-coin interactive  \DIFaddbegin \DIFadd{protocol.  The }\DIFaddend root \(\rt_{ABC}\)
 \DIFaddbegin \DIFadd{fixes three arbitrary interleaved words before the scalar challenges are
sampled; we do not assume that these words are valid encodings.  After
\((r,s)\)}\DIFaddend , the prover  \DIFaddbegin \DIFadd{fixes }\DIFaddend the intermediate-vector  \DIFaddbegin \DIFadd{root before the query set
}\DIFaddend \(I\) and the code-check points  \DIFaddbegin \DIFadd{are sampled}\DIFaddend .

 \DIFaddbegin \DIFadd{For \(W\in\mathbb{F}^{k\times n}\), let
\(\Delta(W,\mathcal{C}^{\langle k\rangle})\) be its relative column distance
from the interleaved code.  We use the following code property; the full
parameter discussion is given in Appendix~\ref{app:security-proof}.
}

\begin{definition}[Structured-fold proximity gap]
\label{def:structured-proximity-gap}
\DIFadd{For \(\vect{a}=(1,a,\ldots,a^{k-1})\), define
\(\mathsf{FoldWord}(\vect{a},W)[j]=\fold(\vect{a},W[*,j])\).  The code has a
\((k,\tau,\epsilon_{\mathsf{gap}})\) structured-fold proximity gap if every
\(W\) at distance greater than \(\tau\) folds to a word within distance
\(\tau\) of \(\mathcal{C}\) }\DIFaddend with probability at most
 \DIFaddbegin \DIFadd{\(\epsilon_{\mathsf{gap}}\) over \(a\sample\mathbb{F}\).
}\end{definition}
\DIFaddend 

 \DIFaddbegin \begin{definition}[Backend assumptions]
\label{def:lamp-backend-assumptions}
\DIFadd{The CP-SNARK, sampled openings, and CP-Link are sound with errors
\(\epsilon_{\mathsf{cp}}\), \(\epsilon_{\mathsf{open}}\), and
\(\epsilon_{\mathsf{link}}\), respectively.  The four intermediate-vector
encoding checks have error \(\epsilon_{\mathsf{code}}\); for }\DIFaddend the
Reed--Solomon  \DIFaddbegin \DIFadd{check of }\DIFaddend Appendix~\ref{app:rs-barycentric} \DIFaddbegin \DIFadd{,
}\[
  \DIFadd{\epsilon_{\mathsf{code}}
  \le \frac{4(n-1)}{|\mathbb{F}|-n}.
}\]
\DIFadd{For zero knowledge, the CP-SNARK and CP-Link are zero knowledge and the
Pedersen commitments are hiding.
}\DIFaddend \end{definition}

\begin{theorem}[Completeness, soundness, and zero knowledge]
\label{thm:security}
 \DIFaddbegin \DIFadd{Let \(\mathcal{C}\) have relative distance \(\delta\), fix
\(0<\tau<\delta/2\), and assume Definition~\ref{def:structured-proximity-gap}
and Definition~\ref{def:lamp-backend-assumptions}.  If \(\rt_{ABC}\) is fixed
before \((r,s)\), then the interactive }\DIFaddend \(\protocol\)  \DIFaddbegin \DIFadd{protocol has perfect
completeness and zero knowledge.  Moreover, for any PPT prover, let
\(\mathsf{BadInput}\) denote the event that an input word is farther than
\(\tau\) from the interleaved code or that the uniquely decoded matrices
satisfy \(\mat{A}\mat{B}\ne\mat{C}\).  Then
}\[
\DIFadd{\begin{aligned}
\Pr[\mathsf{Acc}\land\mathsf{BadInput}]
&\le \epsilon_{\mathsf{cp}}+\epsilon_{\mathsf{open}}
   +\epsilon_{\mathsf{link}}\\
&\quad +\epsilon_{\mathsf{code}}+3\epsilon_{\mathsf{gap}}
   +3(1-\tau)^t\\
&\quad +\frac{k-1}{|\mathbb{F}|}
   +4(1-\delta+\tau)^t .
\end{aligned}
}\]\DIFaddend 
 \end{theorem}

 \begin{proof}[Proof sketch]
 \DIFaddbegin \DIFadd{Condition on the soundness }\DIFaddend of the CP-SNARK,  \DIFaddbegin \DIFadd{openings}\DIFaddend , CP-Link, and
 \DIFaddbegin \DIFadd{intermediate encoding checks.  The roots then bind all words before \(I\), and
the circuit uses their authenticated sampled symbols.
}

\DIFadd{The \(r\)-folds test the input words for \(\mat{A}\) and \(\mat{C}\).  The
coefficient vector \(\vect{x}=\vect{r}\mat{A}\) need not be uniform, so the
\(\vect{x}\)-weighted product equation alone does not give the same guarantee
for \(\mat{B}\); the independent \(s\)-fold supplies that test.  By the
structured-fold gap, any input farther than \(\tau\) from the code is rejected
except for \(3\epsilon_{\mathsf{gap}}+3(1-\tau)^t\)}\DIFaddend .

 \DIFaddbegin \DIFadd{Otherwise, \(\tau<\delta/2\) gives unique decoded matrices.  If their product
relation is false, structured Freivalds fails }\DIFaddend with probability at most
\((k-1)/|\mathbb{F}|\) \DIFaddbegin \DIFadd{.  Outside that event}\DIFaddend , at least one  \DIFaddbegin \DIFadd{sampled relation
disagrees on a \((\delta-\tau)\)}\DIFaddend -fraction of positions,  \DIFaddbegin \DIFadd{contributing
\(4(1-\delta+\tau)^t\).  }\DIFaddend Zero knowledge follows  \DIFaddbegin \DIFadd{from the simulators of the
CP-SNARK and }\DIFaddend CP-Link  \DIFaddbegin \DIFadd{and the hiding of Pedersen commitments.  The complete
bad-event decomposition appears in Appendix~\ref{app:security-proof}.
}\DIFaddend \end{proof}

\begin{remark}[Fiat--Shamir]
The  \DIFaddbegin \DIFadd{implementation derives \((r,s)\)}\DIFaddend , \(I\), and the code-check points from
domain-separated transcript hashes \DIFaddbegin \DIFadd{.  A formal analysis of this multi-round
}\DIFaddend Fiat--Shamir transform  \DIFaddbegin \DIFadd{is outside the interactive theorem and left to future
work}\DIFaddend .
\end{remark}


%% Section 7: Implementation and Evaluation 
\section{Implementation and Evaluation}
\label{sec:implementation}

We implement \(\protocol\) and evaluate its performance. The implementation follows the construction in Section~\ref{sec:construction}.  We instantiate the CP-SNARK with Groth16, using the commit-and-prove
interface~\cite{gnark,legosnark}.  For CP-Link, we use a pairing-based QA-NIZK
proof~\cite{qa-nizk}; the concrete instantiation used in our implementation is
given in Appendix~\ref{app:cplink-instant}.  The implementation is available at
\url{https://anonymous.4open.science/r/LAMP-38E3/.}

Our evaluation is organized around four questions.  First, does the number of constraints of \(\protocol\) scale linearly with the matrix dimension, as predicted by Section~\ref{sec:protocol-complexity}?  Second, which components dominate proving time, proof size, and verification time?  Third, how does \(\protocol\) behave when proving multiple independent claims \(\mat{A}_i\mat{B}_i=\mat{C}_i\) in a batch, rather than a single product \(\mat{A}\mat{B}=\mat{C}\)?  Fourth, does the advantage remain useful on a neural-network workload? \\

\noindent\textbf{Experimental setup.}
We implement both \(\protocol\) and the Freivalds baseline in Go using gnark, with BN254 as the elliptic curve.  All experiments were run on a Linux server with an AMD EPYC 7B13 CPU with 32 cores and 254GB of memory.  We use a Reed--Solomon code of rate \(\rho=1/2\) and set the number of sampled positions to \(t=128\). \DIFaddbegin \\
\DIFaddend 

\DIFaddbegin \DIFadd{Groth16 requires a circuit-specific setup, while the QALink instantiation uses
a configuration-dependent CRS determined by parameters including \(t\) and the
linked block length. These setup procedures are offline preprocessing and are
not included in the online proving and verification times reported below; when
setup time is reported, we list it separately. Changing the circuit or these
QALink parameters requires regenerating the corresponding setup material.}\\




\DIFaddend \begin{table*}[!t]
\centering
\setlength{\tabcolsep}{3.5pt}
\begin{adjustbox}{max width=\textwidth}
\begin{tabular}{c cc cc cc cc}
\toprule
\multirow{2}{*}{\(\boldsymbol{k}\)}
& \multicolumn{2}{c}{\textbf{Constraints}}
& \multicolumn{2}{c}{\textbf{Prove Time}}
& \multicolumn{2}{c}{\textbf{Verify Time}}
& \multicolumn{2}{c}{\textbf{Proof Size}} \\
\cmidrule(lr){2-3}\cmidrule(lr){4-5}\cmidrule(lr){6-7}\cmidrule(lr){8-9}
& \(\protocol\) & Freivalds
& \(\protocol\) & Freivalds
& \(\protocol\) & Freivalds
& \(\protocol\) & Freivalds \\
\midrule
$2^7$    & 54{,}069      & 49{,}610       & 0.76\,s   & 0.58\,s   & 0.06\,s & 1\,ms & 22{,}852\,B & \multirow{6}{*}{196\,B} \\
$2^8$    & 105{,}915     & 197{,}194      & 1.50\,s   & 2.03\,s   & 0.06\,s & 1\,ms & 29{,}380\,B & \\
$2^9$    & 209{,}995     & 788{,}810      & 3.08\,s   & 6.89\,s   & 0.06\,s & 1\,ms & 36{,}292\,B & \\
$2^{10}$ & 418{,}149     & 3{,}156{,}298  & 7.11\,s   & 25.63\,s  & 0.06\,s & 2\,ms & 43{,}332\,B & \\
$2^{11}$ & 834{,}075     & 12{,}622{,}154 & 16.80\,s  & 102.07\,s & 0.06\,s & 2\,ms & 51{,}076\,B & \\
$2^{12}$ & 1{,}666{,}315 & 50{,}495{,}818 & 48.06\,s  & 405.04\,s & 0.06\,s & 2\,ms & 58{,}564\,B & \\
$2^{13}$ & 3{,}330{,}789 & --             & 150.24\,s & --        & 0.06\,s & --    & 66{,}756\,B & -- \\
\bottomrule
\end{tabular}
\end{adjustbox}
\vspace{6pt}
\caption{\(\protocol\) compared with a Freivalds baseline.}
\label{tab:lamp-vs-freivalds}
\end{table*}


\subsection{Comparison with a Freivalds Baseline}
\label{sec:eval-freivalds}
We compare \(\protocol\) with a Freivalds-based baseline implemented as a
SNARK circuit. Table~\ref{tab:lamp-vs-freivalds} reports the number of
constraints, proving time, verification time, and proof size for both systems.
\DIFaddbegin \DIFadd{This baseline is not state of the art; because it uses the same Groth16 stack,
it isolates the effect of moving the Freivalds vector--matrix products out of
the circuit. The Freivalds proof remains \(196\) B at every dimension because a
Groth16 proof is constant-size and does not grow with the number of circuit
constraints. In contrast, a \(\protocol\) proof additionally contains sampled
Merkle openings and CP-Link proofs.
}\DIFaddend 

For small dimensions, the fixed cost of
\(\relation_{\mathsf{LAMP}}^{\mathsf{on}}\) dominates. In particular, at
\(k=2^7\), \(\protocol\) uses slightly more constraints than the baseline.
Starting from \(k=2^8\), however, the asymptotic advantage of \(\protocol\)
becomes visible. As \(k\) doubles, the Freivalds baseline grows by roughly
\(4\times\), whereas \(\protocol\) grows by roughly \(2\times\). This behavior
matches the expected gap between the \(\mathcal{O}(k^2)\) vector--matrix
computation in the baseline and the \(\mathcal{O}(tk)\) sampled fold check in
\(\protocol\).
At \(k=4{,}096\), \(\protocol\) reduces the number of constraints from
\(50{,}495{,}818\) to \(1{,}666{,}315\), corresponding to a \(30.30\times\)
reduction. The proving time also decreases from \(405.04\) s to \(48.06\) s.
At \(k=8{,}192\), the Freivalds baseline runs out of memory, while
\(\protocol\) still completes with \(3{,}330{,}789\) constraints.

The main cost of \(\protocol\) appears in verification and proof size.
Verification takes about \(0.06\) s for \(\protocol\), compared with
\(1\)--\(2\) ms for the Freivalds baseline. This additional overhead comes from
verifying Merkle openings and the pairing-based CP-Link proofs in addition to
the Groth16 proof. In the QA-Link instantiation, CP-Link verification is the
dominant part of this overhead because the verifier evaluates the sampled link
equations as a multi-pairing check. Across all measured matrix dimensions, verification stays at about \(0.06\) s, while the prover-side savings become substantial for large matrix dimensions.


\begin{table*}[!t]
\centering
% \scriptsize
\setlength{\tabcolsep}{3.5pt}
\begin{adjustbox}{max width=\textwidth}
\begin{tabular}{c cc ccc ccc ccc}
\toprule
\multirow{2}{*}{\(\boldsymbol{k}\)}
& \multicolumn{2}{c}{\textbf{Commit Time}}
& \multicolumn{3}{c}{\textbf{Prove Time}}
& \multicolumn{3}{c}{\textbf{Verify Time}}
& \multicolumn{3}{c}{\textbf{Proof Size}} \\
\cmidrule(lr){2-3}\cmidrule(lr){4-6}\cmidrule(lr){7-9}\cmidrule(lr){10-12}
& ABC & XYZ
& Merkle & Groth16 & CP-Link
& Groth16 & Merkle & CP-Link
& Merkle & Groth16 & CP-Link \\
\midrule
$2^7$    & 0.09\,s    & 0.01\,s & 1\,ms & 0.63\,s  & 0.02\,s & 2\,ms & 1\,ms & 0.05\,s & 22{,}464\,B & \multirow{7}{*}{260\,B} & \multirow{7}{*}{128\,B} \\
$2^8$    & 0.25\,s    & 0.02\,s & 1\,ms & 1.18\,s  & 0.05\,s & 2\,ms & 1\,ms & 0.05\,s & 28{,}992\,B & & \\
$2^9$    & 0.87\,s    & 0.03\,s & 1\,ms & 2.11\,s  & 0.07\,s & 2\,ms & 1\,ms & 0.06\,s & 35{,}904\,B & & \\
$2^{10}$ & 3.09\,s    & 0.13\,s & 1\,ms & 3.78\,s  & 0.11\,s & 2\,ms & 1\,ms & 0.05\,s & 42{,}944\,B & & \\
$2^{11}$ & 9.15\,s    & 0.45\,s & 1\,ms & 7.03\,s  & 0.17\,s & 2\,ms & 1\,ms & 0.05\,s & 50{,}688\,B & & \\
$2^{12}$ & 32.32\,s   & 1.81\,s & 1\,ms & 13.60\,s & 0.33\,s & 2\,ms & 2\,ms & 0.05\,s & 58{,}176\,B & & \\
$2^{13}$ & 116.78\,s  & 6.23\,s & 1\,ms & 26.51\,s & 0.72\,s & 2\,ms & 2\,ms & 0.05\,s & 66{,}368\,B & & \\
\bottomrule
\end{tabular}
\end{adjustbox}
\vspace{6pt}
\caption{Component-level measurements for \(\protocol\).}
\label{tab:lamp-breakdown}
\end{table*}

\subsection{LAMP Component Breakdown}
\label{sec:eval-lamp-breakdown}

Table~\ref{tab:lamp-breakdown} breaks down the measured cost of \(\protocol\)
by component.  Commit time includes  \DIFaddbegin \DIFadd{encoding, }\DIFaddend Pedersen leaf commitments\DIFaddbegin \DIFadd{, }\DIFaddend and
Merkle root generation.  For small dimensions, SNARK proving is the main bottleneck.
As \(k\) grows, however, commitment generation for the ABC and XYZ objects
becomes dominant.  For example, at \(k=4{,}096\), commitment generation takes
\(34.13\) s in total, while Groth16 proving takes \(13.60\) s and CP-Link
proving takes \(0.33\) s.  On the verifier side, Groth16 and Merkle
verification remain at a few milliseconds, while the pairing-based QA-Link
check takes about \(0.05\)--\(0.06\) s.  Thus total verification remains about \(0.06\) s for all measured matrix dimensions.



\begin{table*}[!t]
\centering
\scriptsize
\setlength{\tabcolsep}{2.4pt}
\begin{adjustbox}{max width=\textwidth}
\begin{tabular}{c c cc ccc ccc}
\toprule
\multirow{2}{*}{\textbf{Claims}}
& \multirow{2}{*}{\textbf{Constraints}}
& \multicolumn{2}{c}{\textbf{Commit Time}}
& \multicolumn{3}{c}{\textbf{Prove Time}}
& \multicolumn{3}{c}{\textbf{Proof Size}} \\
\cmidrule(lr){3-4}\cmidrule(lr){5-7}\cmidrule(lr){8-10}
& & ABC & XYZ
& Merkle & Groth16 & CP-Link
& Merkle & Groth16 & CP-Link \\
\midrule
1  & 54{,}069  & 0.11\,s & 0.01\,s & 1\,ms & 0.65\,s & 0.02\,s & 22{,}592\,B & \multirow{10}{*}{260\,B} & \multirow{10}{*}{128\,B} \\
2  & 105{,}349 & 0.15\,s & 0.01\,s & 1\,ms & 1.18\,s & 0.04\,s & 22{,}528\,B & & \\
3  & 156{,}629 & 0.19\,s & 0.02\,s & 1\,ms & 1.69\,s & 0.04\,s & 22{,}336\,B & & \\
4  & 207{,}909 & 0.23\,s & 0.01\,s & 1\,ms & 2.14\,s & 0.07\,s & 21{,}696\,B & & \\
5  & 259{,}189 & 0.30\,s & 0.02\,s & 1\,ms & 2.48\,s & 0.08\,s & 22{,}528\,B & & \\
6  & 310{,}469 & 0.32\,s & 0.02\,s & 1\,ms & 2.96\,s & 0.08\,s & 22{,}272\,B & & \\
7  & 361{,}749 & 0.34\,s & 0.02\,s & 1\,ms & 3.35\,s & 0.10\,s & 22{,}208\,B & & \\
8  & 413{,}029 & 0.39\,s & 0.02\,s & 1\,ms & 3.75\,s & 0.12\,s & 22{,}272\,B & & \\
9  & 464{,}309 & 0.42\,s & 0.02\,s & 1\,ms & 4.16\,s & 0.11\,s & 22{,}464\,B & & \\
10 & 515{,}589 & 0.46\,s & 0.02\,s & 1\,ms & 4.60\,s & 0.12\,s & 22{,}208\,B & & \\
\bottomrule
\end{tabular}
\end{adjustbox}
\vspace{6pt}
\caption{Batching \(q\) independent matrix-multiplication claims
\(\{\mat{A}_i\mat{B}_i=\mat{C}_i\}_{i\in[q]}\) at fixed \(k=2^7\).}
\label{tab:lamp-batch}
\end{table*}

\subsection{Batching Matrix-Multiplication Claims}
\label{sec:eval-batch}
We also evaluate a batched setting with multiple independent claims \(\mat{A}_i\mat{B}_i=\mat{C}_i\).  The number of claims is the number of matrix products proved together.  For each codeword position, the prover packs all interleaved symbols of \(\mat{A}_i,\mat{B}_i,\mat{C}_i\) into a single Pedersen leaf and builds one Merkle tree.  Therefore, the number of Merkle membership proofs does not grow with the number of claims.

Table~\ref{tab:lamp-batch} reports the component costs at fixed \(k=2^7\).
As expected, the constraint count and Groth16 proving time grow with the number
of checked matrix products.  In contrast, the Merkle proof size stays around
\(22\) KB across all claim counts.  This shows that, once the batched matrix
leaves are committed together, the sampled Merkle membership proof does not
scale with the number of matrix products.  The small variation in Merkle proof
size comes from the multi-opening implementation: the number of shared
authentication nodes can vary slightly depending on the sampled index set.
Verification time remains about \(0.06\)--\(0.07\) s across the batch-size
range, because the number of sampled positions and the verification procedure
are unchanged.

\subsection{GPT-2 Workload}
\label{sec:eval-gpt2}

Table~\ref{tab:lamp-gpt2} evaluates \(\protocol\) as a verifiable AI
application: the matrix multiplications in a single GPT-2 layer with sequence
length \(2^{10}\).  This workload consists of several matrix multiplications
arising from the attention and MLP components of the layer, with different
matrix shapes and structures.
\(\protocol\) reduces the constraint count from \(76{,}807{,}671\) to
\(33{,}427{,}908\), a \(2.30\times\) reduction, and reduces proving time from
\(408.69\) s to \(230.82\) s, a \(1.77\times\) improvement.  The gain is smaller
than in the square-matrix benchmark because a GPT-2 layer decomposes into
several rectangular matrix multiplications and smaller matrix products inside
each attention head.  Since these products have different dimensions, the
sampled opening layer cannot easily amortize all products through a single
shared leaf layout.
This effect also appears in verifier-side costs.  The Freivalds baseline
verifies in \(0.49\) s, whereas \(\protocol\) uses more commitments and therefore
requires the verifier to check more CP-Link proofs and Merkle membership proofs.
As a result, verification takes \(6.52\) s for \(\protocol\), and the proof size
increases \DIFaddbegin \DIFadd{from \(196\) B }\DIFaddend to \(3{,}342{,}724\) B \DIFaddbegin \DIFadd{(approximately \(3.34\) MB)}\DIFaddend .
Thus, on this heterogeneous AI workload,
\(\protocol\) still reduces prover-side cost, while the sampled opening and
linking layers become the main contributors to verification time and proof size.

\begin{table}[H]
\centering
% \footnotesize
\setlength{\tabcolsep}{3pt}
\begin{tabular*}{\columnwidth}{@{\extracolsep{\fill}}ccccc@{}}
\toprule
\textbf{System} & \textbf{Constraints} & \textbf{Prove} & \textbf{Verify} & \textbf{Proof Size} \\
\midrule
\(\protocol\) & 33{,}427{,}908 & 230.82\,s & 6.52\,s & 3{,}342{,}724\,B \\
Freivalds & 76{,}807{,}671 & 408.69\,s & 0.49\,s & 196\,B \\
\bottomrule
\end{tabular*}
\vspace{6pt}
\caption{GPT-2 at sequence length \(2^{10}\).}
\label{tab:lamp-gpt2}
\end{table}


%% Section 8: Conclusion
\section{Conclusion}
\label{sec:conclusion}

We proposed \(\protocol\), a protocol for verifying matrix multiplication by
moving expensive vector--matrix products outside the SNARK circuit and checking
them through sampled local tests over encoded data.  The prover commits to the
encoded matrices and Freivalds intermediate vectors, while the SNARK circuit
proves only the sampled fold relations.  Merkle openings and CP-Link proofs
ensure that the values used inside the CP-SNARK are consistent with the
externally committed data fixed before the sampling challenges are known.

Our implementation confirms the intended tradeoff.  For a fixed number \(t\) of
queries, the main arithmetic relation has
\(\mathcal{O}(tk+E_{\mathsf{code}})\) constraints rather than the
\(\mathcal{O}(k^2)\) constraints of a Freivalds baseline.  Thus, when \(t\) is
fixed and the code-consistency checks are linear, the constraint count of
\(\protocol\) grows linearly with the matrix dimension \(k\).  In the matrix
benchmark at \(k=4{,}096\), \(\protocol\) reduces the constraint count by
\(30.3\times\) and proving time by \(8.43\times\), while keeping verification
around \(0.06\) s.  The batching experiment further shows that opening overhead
can be amortized when multiple matrix-multiplication claims share the same leaf
layout.

 \DIFaddbegin \DIFadd{On the heterogeneous }\DIFaddend GPT-2  \DIFaddbegin \DIFadd{workload}\DIFaddend , \(\protocol\)  \DIFaddbegin \DIFadd{reduces constraints by
}\DIFaddend \(2.30\times\)  \DIFaddbegin \DIFadd{and proving time by }\DIFaddend \(1.77\times\) \DIFaddbegin \DIFadd{, but verification increases
from \(0.49\) s to \(6.52\) s and proof size from \(196\) B to \(3.34\) MB.
Thus, this experiment demonstrates a prover-side trade-off rather than an
end-to-end performance improvement}\DIFaddend .

The main remaining costs are sampled-opening proof size and pairing-based link
verification.  As future work, we plan to reduce opening and linking overhead,
improve sharing across consecutive matrix products of different sizes, and
exploit fixed model weights in verifiable AI scenarios for preprocessing and
batching.  These optimizations may reduce proof size and verification time on
more complex AI workloads while preserving the linear constraint growth of
\(\protocol\).


\newpage


\bibliographystyle{IEEEtran}
\bibliography{Styles/bibliography}

\appendices
%DIF >  Standard SNARK and CP-Link definitions were included in an earlier draft.
%DIF >  They are omitted from the revised paper because Section~\ref{sec:security}
%DIF >  states exactly the backend properties used by the theorem.
\DIFaddbegin \iffalse
\DIFaddend \section{SNARK Security Definitions}
\label{app:snark-security}

Let \(\relation=\{\relation_\lambda\}_{\lambda\in\mathbb{N}}\) be a family of
NP relations.  A SNARK for \(\relation\) is a tuple of algorithms
\(\bPi=(\setup,\prove,\verify)\), as defined in
Section~\ref{sec:prelim-snark}.  We recall the standard security properties
used in this work. \\

\noindent\textbf{Completeness.}
\(\bPi\) satisfies completeness if, for every security parameter \(\lambda\)
and every \((\stmt,\wit)\in\relation_\lambda\),
\[
\begin{aligned}
&\Pr\!\left[
\begin{array}{l}
  \bPi.\verify(\para,\stmt,\pi)=1
  \\[1mm]
  \para\gets\bPi.\setup(1^\lambda,\relation_\lambda),
  \quad
  \pi\gets\bPi.\prove(\para,\stmt,\wit)
\end{array}
\right]
\\
&\hspace{2em}\ge 1-\textsf{negl}(\lambda).
\end{aligned}
\]
If the probability is \(1\), we say that \(\bPi\) has perfect completeness. \\

\noindent\textbf{Knowledge soundness.}
\(\bPi\) satisfies knowledge soundness if, for every \(\ppt\) prover
\(\mathcal{A}\), there exists  \DIFaddbegin \DIFadd{a }\DIFaddend \(\ppt\) extractor
\(\mathsf{Ext}^{\mathcal{A}}\) such that
\[
\begin{aligned}
&\Pr\left[
\begin{array}{l}
  \bPi.\verify(\para,\stmt,\pi)=1 \\
  \quad {}\wedge{}\,(\stmt,\wit)\notin\relation_\lambda
\end{array}
\;:\;
\begin{array}{l}
  \para\gets\bPi.\setup(1^\lambda,\relation_\lambda),\\
  (\stmt,\pi)\gets\mathcal{A}(\para),\\
  \wit\gets\mathsf{Ext}^{\mathcal{A}}(\para,\stmt)
\end{array}
\right] \\
&\le \textsf{negl}(\lambda).
\end{aligned}
\]
Equivalently, except with negligible probability, any prover that outputs an
accepting proof for \(\stmt\) must know a witness \(\wit\) such that
\((\stmt,\wit)\in\relation_\lambda\). \\


\noindent\textbf{Zero knowledge.}
\(\bPi\) satisfies zero knowledge if there exists a \(\ppt\) simulator
\(\mathsf{Sim}\) such that, for every \(\ppt\) distinguisher
\(\mathcal{D}\), the following two experiments are indistinguishable.

\[
\begin{aligned}
\mathsf{Real}_{\bPi,\mathcal{D}}(\lambda):
\quad
&\left\{
\begin{array}{@{}l@{}}
  \para\gets\bPi.\setup(1^\lambda,\relation_\lambda),\\
  \pi\gets\bPi.\prove(\para,\stmt,\wit),\\
  b\gets\mathcal{D}(\para,\stmt,\pi),\\
  \text{return } b.
\end{array}
\right.
\\[3mm]
\mathsf{Ideal}_{\bPi,\mathcal{D}}(\lambda):
\quad
&\left\{
\begin{array}{@{}l@{}}
  \para\gets\bPi.\setup(1^\lambda,\relation_\lambda),\\
  \pi^\star\gets\mathsf{Sim}(\para,\stmt),\\
  b\gets\mathcal{D}(\para,\stmt,\pi^\star),\\
  \text{return } b.
\end{array}
\right.
\end{aligned}
\]

For all \((\stmt,\wit)\in\relation_\lambda\),
\[
\left|
  \Pr[\mathsf{Real}_{\bPi,\mathcal{D}}(\lambda)=1]
  -
  \Pr[\mathsf{Ideal}_{\bPi,\mathcal{D}}(\lambda)=1]
\right|
\le
\textsf{negl}(\lambda).
\]
A SNARK satisfying zero knowledge is called a zk-SNARK.

\section{CP-Link Security Definitions}
\label{app:cplink-security}

Let
\(\relation_{\mathsf{link}}
=\{\relation_{\mathsf{link},\lambda}\}_{\lambda\in\mathbb{N}}\)
be the blockwise linking relation defined in Section~\ref{sec:prelim-cplink}.
For a SNARK-side key \(\ck_S\), an external key \(\ck_E\), a SNARK-side
commitment \(\widetilde{\cm}\), external commitments
\(\cm_0,\ldots,\cm_{q-1}\), sampled blocks
\((\vect{m}_i)_{i=0}^{q-1}\), and opening randomness
\(\widetilde{o},o_0,\ldots,o_{q-1}\), define
\[
\begin{aligned}
\stmt_{\mathsf{link}}
&=
((\ck_S,\widetilde{\cm}),
(\ck_E,\cm_0,\ldots,\cm_{q-1})),\\
\wit_{\mathsf{link}}
&=
((\vect{m}_i)_{i=0}^{q-1},\widetilde{o},(o_i)_{i=0}^{q-1}).
\end{aligned}
\]
The relation \(\relation_{\mathsf{link},\lambda}\) contains exactly the pairs
\((\stmt_{\mathsf{link}},\wit_{\mathsf{link}})\) satisfying
\[
\begin{aligned}
\widetilde{\cm}
&=
\com(\ck_S,\vect{m}_0\|\cdots\|\vect{m}_{q-1};\widetilde{o}),\\
\cm_i
&=
\com(\ck_E,\vect{m}_i;o_i)
\quad \text{for all } i=0,\ldots,q-1.
\end{aligned}
\]
A CP-Link proof system for \(\relation_{\mathsf{link}}\) is a tuple of
algorithms
\(\bPi_{\mathsf{link}}=(\setup,\prove,\verify)\), as defined in
Section~\ref{sec:prelim-cplink}.  The corresponding linking language is
\[
  \mathcal{L}_{\mathsf{link},\lambda}
  =
  \left\{
    \stmt_{\mathsf{link}}
    :
    \exists \wit_{\mathsf{link}}
    \text{ such that }
    (\stmt_{\mathsf{link}},\wit_{\mathsf{link}})
    \in
    \relation_{\mathsf{link},\lambda}
  \right\}.
\]
Thus a statement is in \(\mathcal{L}_{\mathsf{link},\lambda}\) exactly when
the SNARK-side commitment opens to the ordered concatenation of sampled blocks
and each external commitment opens to the corresponding block.  The main
security property for our protocol is link soundness.  If the CP-Link backend
is implemented as a NIZK or QA-NIZK, as in our implementation, it may
additionally satisfy witness zero knowledge~\cite{legosnark,qa-nizk}. \\

\noindent\textbf{Link completeness.}
\(\bPi_{\mathsf{link}}\) satisfies completeness if, for every security
parameter \(\lambda\) and every
\((\stmt_{\mathsf{link}},\wit_{\mathsf{link}})
\in\relation_{\mathsf{link},\lambda}\),
\[
\begin{aligned}
&\Pr\!\left[
\begin{array}{l}
  \bPi_{\mathsf{link}}.\verify(
    \para_{\mathsf{link}},
    \stmt_{\mathsf{link}},
    \pi_{\mathsf{link}}
  )=1
  \\[1mm]
  \para_{\mathsf{link}}
    \gets
    \bPi_{\mathsf{link}}.\setup(
      1^\lambda,\relation_{\mathsf{link},\lambda}
    ),\\
  \pi_{\mathsf{link}}
    \gets
    \bPi_{\mathsf{link}}.\prove(
      \para_{\mathsf{link}},
      \stmt_{\mathsf{link}},
      \wit_{\mathsf{link}}
    )
\end{array}
\right]
\\
&\hspace{2em}\ge 1-\textsf{negl}(\lambda).
\end{aligned}
\]
If the probability is \(1\), we say that
\(\bPi_{\mathsf{link}}\) has perfect completeness. \\

\noindent\textbf{Link soundness.}
\(\bPi_{\mathsf{link}}\) satisfies link soundness with error
\(\epsilon_{\mathsf{link}}(\lambda)\) if, for every \(\ppt\) adversary
\(\mathcal{A}\),
\[
\begin{aligned}
&\Pr\!\left[
\begin{array}{l}
  \bPi_{\mathsf{link}}.\verify(
    \para_{\mathsf{link}},
    \stmt_{\mathsf{link}}^\star,
    \pi_{\mathsf{link}}^\star
  )=1
  \\[1mm]
  {}\wedge{}\ 
  \stmt_{\mathsf{link}}^\star
  \notin
  \mathcal{L}_{\mathsf{link},\lambda}
  \\[1mm]
  \para_{\mathsf{link}}
    \gets
    \bPi_{\mathsf{link}}.\setup(
      1^\lambda,\relation_{\mathsf{link},\lambda}
    ),\\
  (\stmt_{\mathsf{link}}^\star,\pi_{\mathsf{link}}^\star)
    \gets
    \mathcal{A}(\para_{\mathsf{link}})
\end{array}
\right]
\\
&\hspace{2em}\le \epsilon_{\mathsf{link}}(\lambda).
\end{aligned}
\]
If \(\epsilon_{\mathsf{link}}(\lambda)=\textsf{negl}(\lambda)\), we say that
\(\bPi_{\mathsf{link}}\) is link-sound.  Equivalently, except with probability
\(\epsilon_{\mathsf{link}}(\lambda)\), an accepting CP-Link proof certifies the
existence of blocks \((\vect{m}_i)_{i=0}^{q-1}\), an aggregate opening
\(\widetilde{o}\), and external openings \((o_i)_{i=0}^{q-1}\) such that
\(\widetilde{\cm}\) opens to
\(\vect{m}_0\|\cdots\|\vect{m}_{q-1}\) and each \(\cm_i\) opens to
\(\vect{m}_i\). \\

\noindent\textbf{Witness zero knowledge.}
When privacy of the linked blocks and opening randomness is required, we say
that \(\bPi_{\mathsf{link}}\) satisfies witness zero knowledge if there exists a
\(\ppt\) simulator \(\mathsf{Sim}_{\mathsf{link}}\) such that, for every
\(\ppt\) distinguisher \(\mathcal{D}\), the following two experiments are
indistinguishable.

\[
\begin{aligned}
\mathsf{Real}_{\mathsf{link},\mathcal{D}}(\lambda):
\quad
&\left\{
\begin{array}{@{}l@{}}
  \para_{\mathsf{link}}
    \gets
    \bPi_{\mathsf{link}}.\setup(
      1^\lambda,\relation_{\mathsf{link},\lambda}
    ),\\
  \pi_{\mathsf{link}}
    \gets
    \bPi_{\mathsf{link}}.\prove(
      \para_{\mathsf{link}},
      \stmt_{\mathsf{link}},
      \wit_{\mathsf{link}}
    ),\\
  b\gets
    \mathcal{D}(
      \para_{\mathsf{link}},
      \stmt_{\mathsf{link}},
      \pi_{\mathsf{link}}
    ),\\
  \text{return } b.
\end{array}
\right.
\\[3mm]
\mathsf{Ideal}_{\mathsf{link},\mathcal{D}}(\lambda):
\quad
&\left\{
\begin{array}{@{}l@{}}
  \para_{\mathsf{link}}
    \gets
    \bPi_{\mathsf{link}}.\setup(
      1^\lambda,\relation_{\mathsf{link},\lambda}
    ),\\
  \pi_{\mathsf{link}}^\star
    \gets
    \mathsf{Sim}_{\mathsf{link}}(
      \para_{\mathsf{link}},
      \stmt_{\mathsf{link}}
    ),\\
  b\gets
    \mathcal{D}(
      \para_{\mathsf{link}},
      \stmt_{\mathsf{link}},
      \pi_{\mathsf{link}}^\star
    ),\\
  \text{return } b.
\end{array}
\right.
\end{aligned}
\]

For all
\((\stmt_{\mathsf{link}},\wit_{\mathsf{link}})
\in\relation_{\mathsf{link},\lambda}\),
\[
\left|
  \Pr[\mathsf{Real}_{\mathsf{link},\mathcal{D}}(\lambda)=1]
  -
  \Pr[\mathsf{Ideal}_{\mathsf{link},\mathcal{D}}(\lambda)=1]
\right|
\le
\textsf{negl}(\lambda).
\]
This property says that the linking proof reveals no information about
the blocks \((\vect{m}_i)_{i=0}^{q-1}\) or the openings
\(\widetilde{o},o_0,\ldots,o_{q-1}\), beyond the public fact that the
commitments are blockwise linkable.




\DIFaddbegin \fi
\DIFaddend \section{Detailed Proof of Theorem~\ref{thm:security}}
\label{app:security-proof}

\begin{proof}
We prove the public-coin interactive protocol . The Fiat--Shamir  \DIFaddbegin \DIFadd{transform used
by the implementation is discussed separately }\DIFaddend at the end \DIFaddbegin \DIFadd{and is not covered by
the theorem}\DIFaddend .

\paragraph{Perfect completeness}
Assume that the honest prover holds matrices
\(\mat{A},\mat{B},\mat{C}\in\mathbb{F}^{k\times k}\) with
\(\mat{A}\mat{B}=\mat{C}\).  Let
\(\vect{r}=(1,r,\ldots,r^{k-1})\) \DIFaddbegin \DIFadd{and
\(\vect{s}=(1,s,\ldots,s^{k-1})\)}\DIFaddend , and define
 \DIFaddbegin \[
  \DIFadd{\vect{x}=\vect{r}\mat{A},
  \qquad
  \vect{y}=\vect{x}\mat{B},
  \qquad
  \vect{z}=\vect{r}\mat{C},
  \qquad
  \vect{b}=\vect{s}\mat{B}.
}\]\DIFaddend 
Then \(\vect{y}=\vect{z}\).  The honest prover commits to the packed encoded
matrix blocks \(\{\vect{m}_{ABC,j}\}_{j\in[n]}\) under \(\rt_{ABC}\), commits
to the packed encoded vector blocks  \DIFaddbegin \DIFadd{\(\{\vect{m}_{XYZB,j}\}_{j\in[n]}\) under
\(\rt_{XYZB}\)}\DIFaddend , and supplies the sampled openings and CP-Link proofs required
by Section~\ref{sec:openings-and-cplink}.

For every queried position \(j\in I\), code linearity gives
\[
  \fold(\vect{r},\widehat{\mat{A}}[*,j])
  =
  \widehat{\vect{x}}[j],
\]
and similarly
 \DIFaddbegin \[
\DIFadd{\begin{aligned}
  \fold(\vect{x},\widehat{\mat{B}}[*,j])
  &= \widehat{\vect{y}}[j],\\
  \fold(\vect{r},\widehat{\mat{C}}[*,j])
  &= \widehat{\vect{z}}[j],\\
  \fold(\vect{s},\widehat{\mat{B}}[*,j])
  &= \widehat{\vect{b}}[j].
\end{aligned}
}\]\DIFaddend 
Since \(\vect{y}=\vect{z}\), the equality check
\(\widehat{\vect{y}}[j]=\widehat{\vect{z}}[j]\) also holds.  All SNARK,
Merkle, and CP-Link checks are honestly generated, so the verifier accepts.

\paragraph{Bad-event decomposition}
Let \(\mathsf{Acc}\) be the event that the verifier accepts a false committed
matrix multiplication statement.  We decompose the failure probability into
the following events:
\begin{itemize}
  \item \(\mathsf{Bad}_{\mathsf{SNARK}}\): the CP-SNARK accepts although the
  sampled relation of Section~\ref{sec:snark-relation} is false for the
  committed sampled witnesses;
  \item \(\mathsf{Bad}_{\mathsf{open}}\): a Merkle opening or Pedersen leaf
  opening is accepted in a way that violates position binding or commitment
  binding;
	  \item \(\mathsf{Bad}_{\mathsf{link}}\): a CP-Link proof accepts although the
	  SNARK-side sampled blocks are not the corresponding blocks opened from the
	  Pedersen leaves under \(\rt_{ABC}\) and  \DIFaddbegin \DIFadd{\(\rt_{XYZB}\)}\DIFaddend .
  \item \(\mathsf{Bad}_{\mathsf{code}}\): the probabilistic code-consistency
  checks accept encoded views that are not encodings of their corresponding
  intermediate vectors\DIFaddbegin \DIFadd{;
  }\item \DIFadd{\(\mathsf{Bad}_{\mathsf{gap}}\): an input word farther than
  \(\tau\) from the interleaved code produces a structured fold within
  distance \(\tau\) of \(\mathcal{C}\)}\DIFaddend .
\end{itemize}
By the assumed soundness of the selected backends,
 \DIFaddbegin \[
\DIFadd{\begin{aligned}
\Pr[\mathsf{Bad}_{\mathsf{SNARK}}]
  &\le \epsilon_{\mathsf{cp}}(\lambda),\\
\Pr[\mathsf{Bad}_{\mathsf{open}}]
  &\le \epsilon_{\mathsf{open}}(\lambda),\\
\Pr[\mathsf{Bad}_{\mathsf{link}}]
  &\le \epsilon_{\mathsf{link}}(\lambda),\\
\Pr[\mathsf{Bad}_{\mathsf{code}}]
  &\le \epsilon_{\mathsf{code}}(\lambda),\\
\Pr[\mathsf{Bad}_{\mathsf{gap}}]
  &\le 3\epsilon_{\mathsf{gap}}.
\end{aligned}
}\]\DIFaddend 
Condition on the complement of these events.  Then \(\rt_{ABC}\) fixes a
unique sequence of packed matrix-column commitments at all queried positions,
 \DIFaddbegin \DIFadd{\(\rt_{XYZB}\) }\DIFaddend fixes a unique sequence of packed vector-symbol commitments at
all queried positions, and the sampled values used by the SNARK are exactly the
messages opened from those commitments.  In particular, the adversary cannot
choose different sampled blocks after seeing \(I\) without breaking the
opening or linking layer.

\paragraph{ \DIFaddbegin \DIFadd{Input proximity and extraction}\DIFaddend }
The  root \(\rt_{ABC}\)  \DIFaddbegin \DIFadd{fixes arbitrary words
\(W_A,W_B,W_C\in\mathbb{F}^{k\times n}\) before \((r,s)\) is sampled.
Suppose first that one of them is farther than \(\tau\) from the interleaved
code.  We use }\DIFaddend \(r\)  \DIFaddbegin \DIFadd{for \(W_A,W_C\) and the independent \(s\) for \(W_B\).
Conditioned on \(\neg\mathsf{Bad}_{\mathsf{gap}}\), the corresponding folded
word is farther than \(\tau\) from \(\mathcal{C}\).  The encoded view claimed
by the prover is a codeword conditioned on
\(\neg\mathsf{Bad}_{\mathsf{code}}\), so the two words disagree on more than
\(\tau n\) positions.  Since both are committed before \(I\), the probability
that \(t\) independent queries miss this set is at most \((1-\tau)^t\).
Union-bounding over \(A,B,C\) gives \(3(1-\tau)^t\).
}

\DIFadd{It remains to consider the case in which all three words are within distance
\(\tau\) of the interleaved code.  Because \(\tau<\delta/2\), each is close
to a unique interleaved codeword.  Decode those codewords }\DIFaddend row-wise  \DIFaddbegin \DIFadd{to obtain
fixed }\DIFaddend matrices \(\mat{A},\mat{B},\mat{C}\).  \DIFaddbegin \DIFadd{These are the matrices bound by
the accepted committed statement.
}

\paragraph{\DIFadd{Freivalds randomization}}
\DIFaddend Suppose \(\mat{A}\mat{B}\neq\mat{C}\), and let
\(\mat{D}=\mat{A}\mat{B}-\mat{C}\neq 0\).  Some column
\(\vect{d}\) of \(\mat{D}\) is nonzero.  The event
\(\vect{r}\mat{D}=0\) implies
\[
  P(r)
  \coloneqq
  \langle (1,r,\ldots,r^{k-1}),\vect{d}\rangle
  =
  \sum_{i=0}^{k-1} d_i r^i
  =
  0 .
\]
The polynomial \(P\) is nonzero and has degree at most \(k-1\).  Since \(r\)
is sampled after \(\rt_{ABC}\) is fixed, Schwartz--Zippel gives
\[
  \Pr[\vect{r}\mat{D}=0]
  \le
  \frac{k-1}{|\mathbb{F}|}.
\]

\paragraph{ \DIFaddbegin \DIFadd{Sampled relation checks}\DIFaddend }
Now also condition on \(\vect{r}\mat{D}\neq 0\).  No vectors
\(\vect{x},\vect{y},\vect{z}\) can satisfy all four global relations
\[
  \vect{x}=\vect{r}\mat{A},
  \qquad
  \vect{y}=\vect{x}\mat{B},
  \qquad
  \vect{z}=\vect{r}\mat{C},
  \qquad
  \vect{y}=\vect{z}.
\]
Otherwise they would imply
\(\vect{r}\mat{A}\mat{B}=\vect{r}\mat{C}\), contradicting
\(\vect{r}\mat{D}\neq 0\).

Therefore at least one checked encoded relation is false.  Consider the case
\(\vect{x}\neq\vect{r}\mat{A}\); the other three cases are identical.  The
codewords \(\enc(\vect{x})\) and \(\enc(\vect{r}\mat{A})\) are distinct, so by
relative distance they disagree on at least \(\delta n\) coordinates.  \DIFaddbegin \DIFadd{The
committed word \(W_A\) differs from \(\enc(\mat{A})\) at no more than
\(\tau n\) column positions.  Therefore
\(\mathsf{FoldWord}(\vect{r},W_A)\) and \(\enc(\vect{x})\) disagree at no
fewer than \((\delta-\tau)n\) positions.  }\DIFaddend For any queried coordinate \(j\),
code linearity gives
 \DIFaddbegin \[
  \DIFadd{\enc(\vect{r}\mat{A})}[\DIFadd{j}]
  \DIFadd{=
  \fold(\vect{r},\enc(\mat{A})}[\DIFadd{*,j}]\DIFadd{).
}\]\DIFaddend 
The sampled-opening backend ensures that the SNARK uses the opened
 \DIFaddbegin \DIFadd{\(W_A[*,j]\) }\DIFaddend block under \(\rt_{ABC}\) and the opened
\(\widehat{\vect{x}}[j]\) symbol under  \DIFaddbegin \DIFadd{\(\rt_{XYZB}\)}\DIFaddend .  Thus the fold check
rejects whenever \(j\) lands in the disagreement set.  A uniformly sampled
coordinate misses this set with probability at most  \DIFaddbegin \DIFadd{\(1-\delta+\tau\)}\DIFaddend , and
\(t\) independent samples miss it with probability at most
 \DIFaddbegin \DIFadd{\((1-\delta+\tau)^t\)}\DIFaddend .
Union-bounding over the four possible violated relations gives
 \DIFaddbegin \[
  \DIFadd{4(1-\delta+\tau)^t .
}\]\DIFaddend 

\paragraph{Combining the bounds}
Adding the SNARK, opening, linking, code-consistency, Freivalds, and proximity
failure events gives
 \DIFaddbegin \[
\DIFadd{\begin{aligned}
\Pr[\mathsf{Acc}\land\mathsf{BadInput}]
&\le
\epsilon_{\mathsf{cp}}(\lambda)
+\epsilon_{\mathsf{open}}(\lambda)
+\epsilon_{\mathsf{link}}(\lambda)\\
&\quad+\epsilon_{\mathsf{code}}(\lambda)
+3\epsilon_{\mathsf{gap}}\\
&\quad+3(1-\tau)^t
+\frac{k-1}{|\mathbb{F}|}\\
&\quad+4(1-\delta+\tau)^t .
\end{aligned}
}\]\DIFaddend 
 \DIFaddbegin \DIFadd{The earlier estimate \(4(1-\delta)^t\), which assumes exact input codewords,
does not apply to the revised protocol.  A concrete security level must instead
instantiate \(\tau\), \(\epsilon_{\mathsf{gap}}\), and }\DIFaddend \(t\)  \DIFaddbegin \DIFadd{jointly.  In
particular, the experimental choice \(t=128\) is not claimed here to provide
128-bit soundness until that parameter calculation is completed.
}

\DIFadd{This proves Theorem~\ref{thm:security} for the public-coin interactive protocol}\DIFaddend .

\paragraph{Fiat--Shamir variant}
The  \DIFaddbegin \DIFadd{implementation derives \((r,s)\)}\DIFaddend , \(I\), and backend batch challenges from
domain-separated hashes of the transcript.  \DIFaddbegin \DIFadd{A formal }\DIFaddend random-oracle  \DIFaddbegin \DIFadd{analysis of
this multi-round Fiat--Shamir transform is outside }\DIFaddend Theorem~\ref{thm:security}
\DIFaddbegin \DIFadd{and is left to future work}\DIFaddend .
\end{proof}

\section{Barycentric Reed--Solomon Consistency Check}
\label{app:rs-barycentric}

This appendix specifies the Reed--Solomon instantiation of
\(\mathsf{EncCheck}_{\mathcal{C}}\) used in
Section~\ref{sec:snark-relation}.  Let
% \[
%  D=\{\omega_1,\ldots,\omega_n\}\subset\mathbb{F}
% \]
\[
D=\{\omega_0,\ldots,\omega_{n-1}\}\subset\mathbb{F}
\]
be the public evaluation domain.  A message
\(\vect{v}=(v_0,\ldots,v_{k-1})\in\mathbb{F}^k\) defines the degree-\(<k\)
polynomial
 \DIFaddbegin \[
  \DIFadd{p_{\vect{v}}(X)=\sum_{\ell=0}^{k-1} v_\ell X^\ell,
  \qquad
  \enc(\vect{v})=(p_{\vect{v}}(\omega_i))_{i=0}^{n-1}.
}\]\DIFaddend 
Given a claimed encoded word
\(\widehat{\vect{v}}\in\mathbb{F}^n\), let
\(P_{\widehat{\vect{v}}}\) be the unique degree-\(<n\) polynomial satisfying
\[
  P_{\widehat{\vect{v}}}(\omega_i)=\widehat{\vect{v}}[i]
  \qquad\text{for all }i\in[n].
\]
The check samples \(\alpha\sample\mathbb{F}\setminus D\) after the claimed
word is fixed and compares \(P_{\widehat{\vect{v}}}(\alpha)\) with
\(p_{\vect{v}}(\alpha)\).

For efficient evaluation, define the domain polynomial and barycentric weights
 \DIFaddbegin \[
  \DIFadd{L_D(X)=\prod_{i=0}^{n-1}(X-\omega_i),
  \quad
  \lambda_i
  =
  \left(\prod_{\ell\neq i}(\omega_i-\omega_\ell)\right)^{-1}.
}\]\DIFaddend 
For every \(\alpha\notin D\), barycentric interpolation gives
 \DIFaddbegin \[
  \DIFadd{P_{\widehat{\vect{v}}}(\alpha)
  =
  L_D(\alpha)
  \sum_{i=0}^{n-1}
  \lambda_i
  \frac{\widehat{\vect{v}}[i]}{\alpha-\omega_i}.
}\]\DIFaddend 
Thus the concrete Reed--Solomon consistency check is
 \DIFaddbegin \[
\DIFadd{\begin{aligned}
\mathsf{EncCheck}_{\mathsf{RS}}
(\widehat{\vect{v}},\vect{v};\alpha)=1
\Longleftrightarrow{}&
\sum_{\ell=0}^{k-1} v_\ell \alpha^\ell
\\
&=
L_D(\alpha)
\sum_{i=0}^{n-1}
\lambda_i
\frac{\widehat{\vect{v}}[i]}{\alpha-\omega_i}.
\end{aligned}
}\]\DIFaddend 
The right-hand side is a linear form in the claimed encoded symbols once
\(\alpha\) is fixed.  Equivalently, the verifier can derive the interpolation
coefficients
\[
  \gamma_i(\alpha)
  =
  L_D(\alpha)\lambda_i(\alpha-\omega_i)^{-1}
\]
and the circuit checks
 \DIFaddbegin \[
  \DIFadd{\sum_{\ell=0}^{k-1} v_\ell \alpha^\ell
  =
  \sum_{i=0}^{n-1}\gamma_i(\alpha)\widehat{\vect{v}}[i].
}\]\DIFaddend 
This costs \(\mathcal{O}(k+n)\) field operations per checked point, after the
domain-dependent weights are precomputed.

\begin{lemma}[Barycentric code-consistency soundness]
\label{lem:rs-barycentric-soundness}
If \(\widehat{\vect{v}}\neq\enc(\vect{v})\) and
\(\alpha\sample\mathbb{F}\setminus D\) is sampled after
\(\widehat{\vect{v}}\) and \(\vect{v}\) are fixed, then
\[
  \Pr[
    \mathsf{EncCheck}_{\mathsf{RS}}
    (\widehat{\vect{v}},\vect{v};\alpha)=1
  ]
  \le
  \frac{n-1}{|\mathbb{F}|-n}.
\]
\end{lemma}

\begin{proof}
If \(\widehat{\vect{v}}\neq\enc(\vect{v})\), then the interpolation polynomial
\(P_{\widehat{\vect{v}}}\) differs from the degree-\(<k\) polynomial
\(p_{\vect{v}}\).  Their difference
\[
  Q(X)=P_{\widehat{\vect{v}}}(X)-p_{\vect{v}}(X)
\]
is a nonzero polynomial of degree at most \(n-1\).  The check accepts exactly
when \(Q(\alpha)=0\).  Since \(\alpha\) is uniform over
\(\mathbb{F}\setminus D\), at most \(n-1\) choices of \(\alpha\) can make the
check accept, which proves the bound.
\end{proof}

In \(\protocol\), this check is applied independently to
 \DIFaddbegin \DIFadd{\(\widehat{\vect{x}},\widehat{\vect{y}},\widehat{\vect{z}},
\widehat{\vect{b}}\) at
\(\alpha_x,\alpha_y,\alpha_z,\alpha_b\)}\DIFaddend .  A union bound gives the
 \DIFaddbegin \DIFadd{\(4(n-1)/(|\mathbb{F}|-n)\) }\DIFaddend term used in
Definition~\ref{def:lamp-backend-assumptions}.

\DIFaddbegin \iffalse
\DIFaddend \section{CP-Link Instantiation}
\label{app:cplink-instant}

This appendix describes the CP-Link instantiation used in our implementation.
We use a pairing-based quasi-adaptive link argument, denoted QALink, for
linking one SNARK-side Pedersen commitment to several external Pedersen
commitments.  The verifier can also batch multiple QALink verification
equations into one multi-pairing check.

Let \(e:\mathbb{G}_1\times\mathbb{G}_2\to\mathbb{G}_T\) be a bilinear pairing
over prime-order groups of order \(|\mathbb{F}|\).  We write the group
operation in \(\mathbb{G}_1\) additively, and let \(g_2\in\mathbb{G}_2\) be a
generator. \\

\noindent\textbf{Statement and witness.}
Fix a block count \(q\) and a block length \(\ell\).  Let
\(\ck_S=((G^S_{i,\nu})_{i\in[q],\nu\in[\ell]},H^S)\) be the SNARK-side
commitment key for a commitment to \(q\) concatenated blocks, and let
\(\ck_E=((G^E_{\nu})_{\nu\in[\ell]},H^E)\) be the external Pedersen commitment
key for one block.

The public statement is
\(\stmt_{\mathsf{link}}=((\ck_S,\widetilde{\cm}),
(\ck_E,\cm_0,\ldots,\cm_{q-1}))\).  The witness is
\[
  \wit_{\mathsf{link}}
  =
  \big(
    \vect{m}_0,\ldots,\vect{m}_{q-1},
    \widetilde{o},
    o_0,\ldots,o_{q-1}
  \big),
\]
where each block \(\vect{m}_i\in\mathbb{F}^{\ell}\).  The intended relation is
that \(\widetilde{\cm}\) commits to the ordered concatenation
\(\vect{m}_0\|\cdots\|\vect{m}_{q-1}\) under \(\ck_S\), while each
\(\cm_i\) commits to the corresponding block \(\vect{m}_i\) under \(\ck_E\).
Equivalently,
\[
\begin{aligned}
  \widetilde{\cm}
  &=
  \sum_{i\in[q]}\sum_{\nu\in[\ell]}
  \vect{m}_i[\nu]G^S_{i,\nu}
  +
  \widetilde{o}H^S, \\
  \cm_i
  &=
  \sum_{\nu\in[\ell]}
  \vect{m}_i[\nu]G^E_{\nu}
  +
  o_iH^E
  \qquad \text{for every } i\in[q].
\end{aligned}
\]

\noindent\textbf{Setup.}
The setup algorithm takes \((\ck_S,\ck_E,q,\ell)\) as input.  It samples
\(k_0,k_1,\ldots,k_q,a\sample\mathbb{F}^{\ast}\) and defines
\[
  P_{i,\nu} \coloneqq k_0G^S_{i,\nu}+k_{i+1}G^E_{\nu}
  \qquad
  (i\in[q],\,\nu\in[\ell]).
\]
It outputs
\[
\begin{aligned}
  \pk_{\mathsf{link}}
  &=
  \Big(
    (P_{i,\nu})_{i\in[q],\,\nu\in[\ell]},
    k_0H^S,
    (k_{i+1}H^E)_{i\in[q]}
  \Big), \\
  \vk_{\mathsf{link}}
  &=
  (C_0,C_1,\ldots,C_q,A)
  =
  (ak_0g_2,ak_1g_2,\ldots,ak_qg_2,ag_2).
\end{aligned}
\]
The trapdoors \(k_0,\ldots,k_q,a\) are discarded after setup. \\

\noindent\textbf{Proving.}
Given \(\stmt_{\mathsf{link}}\) and \(\wit_{\mathsf{link}}\), the prover
outputs one group element
\[
\begin{aligned}
  \pi_{\mathsf{link}}
  \coloneqq
  \sum_{i\in[q]}\sum_{\nu\in[\ell]}
  \vect{m}_i[\nu]P_{i,\nu}
  +
  \widetilde{o}(k_0H^S)
  +
  \sum_{i\in[q]} o_i(k_{i+1}H^E).
\end{aligned}
\]
By construction and by the homomorphism of Pedersen commitments,
\[
  \pi_{\mathsf{link}}
  =
  k_0\widetilde{\cm}
  +
  \sum_{i\in[q]} k_{i+1}\cm_i .
\]
This identity is the algebraic relation checked by the verifier. \\

\noindent\textbf{Verification.}
The verifier accepts if and only if
\[
  e(\widetilde{\cm},C_0)
  \cdot
  \prod_{i\in[q]} e(\cm_i,C_{i+1})
  =
  e(\pi_{\mathsf{link}},A).
\]
Equivalently, the verifier checks
\[
  e(\widetilde{\cm},C_0)
  \cdot
  \prod_{i\in[q]} e(\cm_i,C_{i+1})
  \cdot
  e(\pi_{\mathsf{link}},-A)
  =
  1_{\mathbb{G}_T}.
\]

\noindent\textbf{Correctness.}
For an honestly generated proof, we have
\(\pi_{\mathsf{link}}=k_0\widetilde{\cm}+\sum_{i\in[q]}k_{i+1}\cm_i\).
Therefore, by bilinearity,
\[
\begin{aligned}
  e(\pi_{\mathsf{link}},A)
  &=
  e\!\left(k_0\widetilde{\cm}+\sum_{i\in[q]}k_{i+1}\cm_i,ag_2\right)
  \\
  &=
  e(\widetilde{\cm},ak_0g_2)
  \cdot
  \prod_{i\in[q]} e(\cm_i,ak_{i+1}g_2)
  \\
  &=
  e(\widetilde{\cm},C_0)
  \cdot
  \prod_{i\in[q]} e(\cm_i,C_{i+1}).
\end{aligned}
\]
Thus the verifier accepts. \\

\noindent\textbf{Soundness.}
QALink is a quasi-adaptive linear proof for the subspace relation induced by
the commitment keys \((\ck_S,\ck_E)\), block count \(q\), and block length
\(\ell\).  Under the standard quasi-adaptive linear-algebra assumption for the
above pairing CRS, any accepting proof for a statement outside this relation
implies a nontrivial linear relation among the hidden setup trapdoors
\(k_0,\ldots,k_q,a\).  Thus the argument is sound except with negligible
probability in the group order. \\

\noindent\textbf{Verifier cost.}
For one link statement with \(q\) external commitments, verification consists
of \(q+2\) pairing terms:
one term for \(\widetilde{\cm}\), \(q\) terms for
\(\cm_0,\ldots,\cm_{q-1}\), and one term for \(\pi_{\mathsf{link}}\) against
\(-A\).  A multi-pairing implementation evaluates these terms with \(q+2\)
Miller-loop inputs and one final exponentiation.  The proof contains one
\(\mathbb{G}_1\) element, and the verifying key contains \(q+2\)
\(\mathbb{G}_2\) elements. \\

\noindent\textbf{Batching pairing equations.}
Suppose the verifier must check \(s\) independent QALink statements.  For the
\(h\)-th statement, let
\[
  E_h
  \coloneqq
  e(\widetilde{\cm}^{(h)},C^{(h)}_0)
  \cdot
  \prod_{i\in[q_h]} e(\cm^{(h)}_i,C^{(h)}_{i+1})
  \cdot
  e(\pi^{(h)}_{\mathsf{link}},-A^{(h)}).
\]
The unbatched verifier checks \(E_h=1_{\mathbb{G}_T}\) for every \(h\).  The
batched verifier first derives nonzero coefficients
\(\lambda_1,\ldots,\lambda_s\in\mathbb{F}^{\ast}\) by hashing all public
statements, verifying keys, proofs, and the surrounding transcript context.  It
then checks the single equation
\[
  \prod_{h=1}^{s} E_h^{\lambda_h}=1_{\mathbb{G}_T}.
\]
Operationally, this is implemented by scaling the \(\mathbb{G}_1\) input of
every pairing term in the \(h\)-th equation by \(\lambda_h\), and then running
one multi-pairing check over all scaled terms:
\[
\begin{aligned}
  &\prod_{h=1}^{s}
  e(\lambda_h\widetilde{\cm}^{(h)},C^{(h)}_0)
  \cdot
  \prod_{h=1}^{s}\prod_{i\in[q_h]}
  e(\lambda_h\cm^{(h)}_i,C^{(h)}_{i+1})
  \cdot \\
  &\prod_{h=1}^{s}
  e(\lambda_h\pi^{(h)}_{\mathsf{link}},-A^{(h)})
  =
  1_{\mathbb{G}_T}.
\end{aligned}
\]

The batched verifier uses \(\sum_{h=1}^{s}(q_h+2)\) Miller-loop terms and one
final exponentiation.  It also performs \(\sum_{h=1}^{s}(q_h+2)\) scalar
multiplications in \(\mathbb{G}_1\) to apply the batching coefficients.  If all
individual equations are valid, then the batched equation is valid by
bilinearity.  If at least one individual equation is invalid, write
\(E_h=g_T^{\delta_h}\) in the prime-order target group.  The batched equation
accepts only if \(\sum_h \lambda_h\delta_h=0\).  Since at least one
\(\delta_h\neq 0\), a fresh random nonzero coefficient vector satisfies this
linear equation with probability at most \(1/(|\mathbb{F}|-1)\).  In the
non-interactive implementation this probability is interpreted in the random
oracle model for the Fiat--Shamir-derived batching coefficients\DIFaddbegin \DIFadd{.
}\fi

\section{\DIFadd{CP-Link Instantiation}}
\label{app:cplink-instant}

\DIFadd{Our implementation uses a pairing-based quasi-adaptive link argument
(QALink).  For \(q\) blocks \(\vect{m}_i\in\mathbb{F}^{\ell}\), the public
statement contains one SNARK-side Pedersen commitment \(\widetilde{\cm}\) and
external commitments \(\cm_0,\ldots,\cm_{q-1}\) satisfying
}\[
\DIFadd{\begin{aligned}
\widetilde{\cm}
 &=\sum_{i\in[q]}\sum_{\nu\in[\ell]}
   \vect{m}_i[\nu]G^S_{i,\nu}+\widetilde{o}H^S,\\
\cm_i
 &=\sum_{\nu\in[\ell]}\vect{m}_i[\nu]G^E_\nu+o_iH^E.
\end{aligned}
}\]
\DIFadd{Setup samples nonzero \(k_0,\ldots,k_q,a\in\mathbb{F}\), publishes the
proving-key elements
\(P_{i,\nu}=k_0G^S_{i,\nu}+k_{i+1}G^E_\nu\), and publishes
\(C_i=ak_ig_2\) and \(A=ag_2\) in the verification key.  The prover outputs
}\[
 \DIFadd{\pi_{\mathsf{link}}
 =k_0\widetilde{\cm}+\sum_{i\in[q]}k_{i+1}\cm_i,
}\]
\DIFadd{expressed using the proving key and the openings.  The verifier checks
}\[
 \DIFadd{e(\widetilde{\cm},C_0)
 \prod_{i\in[q]}e(\cm_i,C_{i+1})
 =e(\pi_{\mathsf{link}},A).
}\]
\DIFadd{Soundness follows from the quasi-adaptive linear-subspace assumption for this
CRS.  One proof contains one \(\mathbb{G}_1\) element; verification uses
\(q+2\) pairing terms, which a multi-pairing implementation evaluates with one
final exponentiation.  The proving and verification keys depend on \(q\) and
\(\ell\), so changing the sample count or linked block layout requires new
setup material.
}

\DIFadd{Multiple link equations are batched by deriving a nonzero transcript
coefficient \(\lambda_h\) for each equation \(E_h=1\) and checking
\(\prod_h E_h^{\lambda_h}=1\) in one multi-pairing.  If at least one equation
is invalid, a fresh coefficient vector satisfies the resulting nontrivial
linear equation with probability at most \(1/(|\mathbb{F}|-1)\).  The
non-interactive implementation derives these coefficients with Fiat--Shamir}\DIFaddend .


\end{document}
